\section{Estudo cohomológico dos morfismos projetivos}
\label{section:III.2}


\subsection{Cálculos explícitos de certos grupos de cohomologia}
\label{subsection:III.2.1}

\begin{env}[2.1.1]
\label{III.2.1.1}
\oldpage[III]{95}
Seja $X$ um pré-esquema e seja $\sh{L}$ um $\sh{O}_X$-módulo invertível; consideremos o anel graduado \sref[0]{0.5.4.6}
\[
\label{III.2.1.1.1}
  S = \Gamma_*(X,\sh{L}) = \bigoplus_{n\in\bb{Z}} \Gamma(X,\sh{L}^{\otimes n}).
  \tag{2.1.1.1}
\]

Seja $(f_i)_{1\leq i\leq r}$ uma família finita de elementos \emph{homogêneos} de $S$, digamos $f_i\in S_{d_i}$; ponhamos $U_i = X_{f_i}$, $U = \bigcup_i U_i$, e denotemos por $\mathfrak{U}$ o recobrimento $(U_i)$ de $U$.
Para todo $\sh{O}_X$-módulo quase-coerente $\sh{F}$, ponhamos
\[
\label{III.2.1.1.2}
  \HH^\bullet(U,\sh{F}(*)) = \bigoplus_{n\in\bb{Z}} \HH^\bullet(U,\sh{F}\otimes\sh{L}^{\otimes n})
  \tag{2.1.1.2}
\]
\[
\label{III.2.1.1.3}
  \HH^\bullet(\mathfrak{U},\sh{F}(*)) = \bigoplus_{n\in\bb{Z}} \HH^\bullet(\mathfrak{U},\sh{F}\otimes\sh{L}^{\otimes n}).
  \tag{2.1.1.3}
\]

Os grupos abelianos (\hyperref[III.2.1.1.2]{2.1.1.2}) e (\hyperref[III.2.1.1.3]{2.1.1.3}) são \emph{bigraduados}, tomando
\[
  (\HH^\bullet(U,\sh{F}(*)))_{mn} = \HH^m(U,\sh{F}\otimes\sh{L}^{\otimes n})
\]
e uma definição análoga para (\hyperref[III.2.1.1.3]{2.1.1.3}).
Para a segunda graduação, é claro que esses grupos são $S$-módulos graduados, como se deduz, por exemplo, do fato de que $\sh{F}\mapsto\HH^m(U,\sh{F})$ e $\sh{F}\mapsto\HH^m(\mathfrak{U},\sh{F})$ são funtores.
\end{env}

\begin{env}[2.1.2]
\label{III.2.1.2}
Consideremos agora o $S$-módulo graduado \sref[0]{0.5.4.6}
\[
\label{III.2.1.2.1}
  M = \Gamma_*(\sh{L},\sh{F}) = \HH^0(X,\sh{F}(*)) = \bigoplus_{n\in\bb{Z}} \Gamma(X,\sh{F}\otimes\sh{L}^{\otimes n}).
  \tag{2.1.2.1}
\]

\oldpage[III]{96}
Se $X$ é um pré-esquema cujo espaço subjacente é noetheriano, ou um esquema quase compacto, deduz-se de \sref[I]{I.9.3.1} que, colocando como de costume $U_{i_0 i_1\cdots i_p} = \bigcap_{k=0}^p U_{i_k}$, tem-se, salvo isomorfismo canônico,
\[
  \Gamma(U_{i_0 i_1\cdots i_p},\sh{F}(*)) = \HH^0(U_{i_0 i_1\cdots i_p},\sh{F}(*)) = M_{f_{i_0}f_{i_1}\cdots f_{i_p}}.
\]

Além disso, com a notação de (\hyperref[III.1.2.2]{1.2.2}), podemos identificar $M_{f_{i_0}f_{i_1}\cdots f_{i_p}}$ com $\varinjlim_n M_{i_0 i_1\cdots i_p}^{(n)}$.
Esta identificação é um isomorfismo de $S$-módulos \emph{graduados} se definirmos o grau de um elemento homogêneo $z\in\varinjlim_n M_{i_0 i_1\cdots i_p}^{(n)}$ da seguinte forma: $z$ é a imagem canônica de um elemento homogêneo $x\in M_{i_0 i_1\cdots i_p}^{(n)} = M$ de grau $m$, e tomamos como grau de $z$ o número $m-n(d_{i_0}+d_{i_1}+\cdots+d_{i_p})$.
Considerando a definição dos homomorfismos $\vphi_{kh}:M_{i_0 i_1\cdots i_p}^{(h)}\to M_{i_0 i_1\cdots i_p}^{(k)}$ (\hyperref[III.1.2.2]{1.2.2}), vê-se imediatamente que esta definição não depende do ``representante'' $x$ de $z$ escolhido.
Denotando, como em (\hyperref[III.1.2.2]{1.2.2}), por $C_n^p(M)$ o conjunto de aplicações alternadas de $[1,r]^{p+1}$ em $M$ (para todo $n$), definimos de forma análoga uma estrutura de $S$-módulo graduado sobre $\varinjlim_n C_n^p(M)$; como em (\hyperref[III.1.2.2]{1.2.2}), temos então
\[
\label{III.2.1.2.2}
  C^p(\mathfrak{U},\sh{F}(*)) = \varinjlim_n C_n^p(M)
  \tag{2.1.2.2}
\]
onde o isomorfismo entre ambos os membros \emph{respeita os graus}.
Tem-se também, como em (\hyperref[III.1.2.2]{1.2.2}),
\[
\label{III.2.1.2.3}
  C^p(\mathfrak{U},\sh{F}(*)) = C^{p+1}((\mathbf{f}),M) = \varinjlim_n K^{p+1}(\mathbf{f}^n,M)
  \tag{2.1.2.3}
\]
onde o isomorfismo \emph{conserva os graus}: o grau de um elemento de $\varinjlim_n K^{p+1}(\mathbf{f}^n,M)$ que seja a imagem canônica de uma cocadeia $g\in K^{p+1}(\mathbf{f}^n,M)$ cujos valores $g(i_0,\dots,i_p)$ pertençam a uma mesma componente homogênea $M_m$ de $M$ é $m-n(d_{i_0}+\cdots+d_{i_p})$, independentemente da escolha dessa cocadeia como representante do elemento em causa.

Como os isomorfismos anteriores são compatíveis com os operadores de coborde, conclui-se, como em (\hyperref[III.1.2.2]{1.2.2}), que:
\end{env}

\begin{proposition}[2.1.3]
\label{III.2.1.3}
Seja $X$ um pré-esquema cujo espaço subjacente é noetheriano, ou um esquema quase compacto.
Existe um isomorfismo canônico, funtorial em $\sh{F}$,
\[
\label{III.2.1.3.1}
  \HH^p(\mathfrak{U},\sh{F}(*)) \simto \HH^{p+1}((\mathbf{f}),M)
  \qquad\text{para todo } p\geq 1.
  \tag{2.1.3.1}
\]
Além disso, existe uma sequência exata, funtorial em $\sh{F}$,
\[
\label{III.2.1.3.2}
  0 \to \HH^0((\mathbf{f}),M) \to M \to \HH^0(\mathfrak{U},\sh{F}(*)) \to \HH^1((\mathbf{f}),M) \to 0.
  \tag{2.1.3.2}
\]
Todos os homomorfismos introduzidos são de grau $0$ para as estruturas de $S$-módulo graduado, onde $S$ é o anel (\hyperref[III.2.1.1.1]{2.1.1.1}).
\end{proposition}

\begin{corollary}[2.1.4]
\label{III.2.1.4}
\oldpage[III]{97}
Se $X$ é um esquema quase-compacto e os $U_i=X_{f_i}$ são afins, existe um isomorfismo canônico, funtorial em $\sh{F}$ e de grau $0$,
\[
\label{III.2.1.4.1}
  \HH^p(U,\sh{F}(*)) \simto \HH^{p+1}((\mathbf{f}),M)
  \qquad\text{para todo } p\geq 1,
  \tag{2.1.4.1}
\]
e uma sequência exata, funtorial em $\sh{F}$,
\[
\label{III.2.1.4.2}
  0 \to \HH^0((\mathbf{f}),M) \to M \to \HH^0(U,\sh{F}(*)) \to \HH^1((\mathbf{f}),M) \to 0,
  \tag{2.1.4.2}
\]
na qual todos os homomorfismos são de grau $0$.
\end{corollary}

\begin{proof}
Basta aplicar \sref{III.1.4.1} ao resultado de \sref{III.2.1.3}.
\end{proof}

A proposição ``local'' análoga a \sref{III.2.1.3} é a seguinte.

\begin{proposition}[2.1.5]
\label{III.2.1.5}
Seja $S$ um anel graduado positivamente, sejam os $f_i$ ($1\leq i\leq r$) elementos homogêneos de $S_+$ de graus $d_i$, e seja $M$ um $S$-módulo graduado.
Seja $X=\Proj(S)$ o espectro primo homogêneo de $S$, e ponhamos $U_i=\mathrm{D}_+(f_i)$, $U=\bigcup_i U_i$ e $\HH^\bullet(U,\widetilde{M}(*))=\bigoplus_{n\in\bb{Z}}\HH^\bullet(U,(M(n))^\sim)$.
Existem isomorfismos canônicos, funtoriais em $M$ e de grau $0$ para estruturas de $S$-módulo graduado,
\[
\label{III.2.1.5.1}
  \HH^p(U,\widetilde{M}(*)) \simto \HH^{p+1}((\mathbf{f}),M)
  \qquad\text{para todo } p\geq 1,
  \tag{2.1.5.1}
\]
e uma sequência exata, funtorial em $M$,
\[
\label{III.2.1.5.2}
  0 \to \HH^0((\mathbf{f}),M) \to M \to \HH^0(U,\widetilde{M}(*)) \to \HH^1((\mathbf{f}),M) \to 0,
  \tag{2.1.5.2}
\]
na qual todos os homomorfismos são de grau $0$.
\end{proposition}

\begin{proof}
De fato, por definição \sref[II]{II.2.5.2},
\[
  \Gamma(U_{i_0i_1\cdots i_p},(M(n))^\sim)
  =(M_{f_{i_0}\cdots f_{i_p}})_n,
\]
e, portanto, $\Gamma(U_{i_0i_1\cdots i_p},\widetilde{M}(*))=M_{f_{i_0}\cdots f_{i_p}}$.
O resto do raciocínio é o mesmo que na demonstração de \sref{III.2.1.4}, levando em conta que $X$ é um esquema.
\end{proof}

\begin{env}[2.1.6]
\label{III.2.1.6}
\emph{Observações.} ---
(i) Sob as hipóteses de \sref{III.2.1.5}, os funtores $\Gamma(U_{i_0i_1\cdots i_p},\widetilde{M}(*))$ são exatos em $M$, por \sref[0]{0.1.3.2}.
O mesmo raciocínio que em (\hyperref[III.1.2.5]{1.2.5}) mostra, portanto, que, se $0\to M'\to M\to M''\to0$ é uma sequência exata de $S$-módulos graduados cujos homomorfismos têm grau $0$, então para todo $p\geq0$ existem diagramas comutativos
\[
\label{III.2.1.6.1}
  \xymatrix{
    \HH^p(U,\widetilde{M''}(*)) \ar[r]^-\partial \ar[d] &
    \HH^{p+1}(U,\widetilde{M'}(*)) \ar[d] \\
    \HH^{p+1}((\mathbf{f}),M'') \ar[r]^-\partial &
    \HH^{p+2}((\mathbf{f}),M')
  }
  \tag{2.1.6.1}
\]

(ii) A Proposição~\sref{III.2.1.5} é especialmente útil quando $S$ é uma $A$-álgebra gerada por um número finito de elementos de grau $1$, com $A$ noetheriano; nesse caso, todo $\sh{O}_X$-módulo quase-coerente é da forma $\widetilde{M}$ \sref[II]{II.2.7.7}.
\end{env}

\begin{env}[2.1.7]
\label{III.2.1.7}
\oldpage[III]{98}
Aplicamos agora \sref{III.2.1.5} ao caso $S=A[\mathrm{T}_0,\dots,\mathrm{T}_r]$, onde $A$ é um anel arbitrário e os $\mathrm{T}_i$ são indeterminados, tomando $M=S$ e $f_i=\mathrm{T}_i$.
Ficamos assim essencialmente reduzidos a calcular $\HH^\bullet((\mathbf{T}),S)$, onde $\mathbf{T}=(\mathrm{T}_i)_{0\leq i\leq r}$.
\end{env}

\begin{lemma}[2.1.8]
\label{III.2.1.8}
Se $S=A[\mathrm{T}_0,\dots,\mathrm{T}_r]$, então, com $\mathbf{T}=(\mathrm{T}_i)_{0\leq i\leq r}$,
\[
\label{III.2.1.8.1}
  \HH^i(\mathbf{T}^n,S)=0
  \qquad\text{se }i\neq r+1,
  \tag{2.1.8.1}
\]
\[
\label{III.2.1.8.2}
  \HH^{r+1}(\mathbf{T}^n,S)=S/(\mathbf{T}^n).
  \tag{2.1.8.2}
\]
Assim, o $A$-módulo $\HH^{r+1}(\mathbf{T}^n,S)$ tem uma $A$-base formada pelas classes módulo $(\mathbf{T}^n)$ dos monômios
\[
  \mathrm{T}_{\mathbf{p}}=\mathrm{T}_0^{p_0}\cdots\mathrm{T}_r^{p_r},
  \qquad \mathbf{p}=(p_0,\dots,p_r),\quad 0\leq p_i<n\ \text{Para tudo }i.
\]
\end{lemma}

\begin{proof}
Isto deduz-se imediatamente de (\hyperref[III.1.1.3.5]{1.1.3.5}) e da Proposição~\sref{III.1.1.4}, cujas hipóteses são satisfeitas trivialmente.
\end{proof}

\begin{env}[2.1.9]
\label{III.2.1.9}
Passando ao limite indutivo em relação a $n$, as relações (\hyperref[III.2.1.8.1]{2.1.8.1}) dão $\HH^i((\mathbf{T}),S)=0$ para $i\neq r+1$.
Para $i=r+1$, o sistema indutivo consiste nos módulos $S/(\mathbf{T}^n)$, sendo o homomorfismo
\[
  \vphi_{nm}:S/(\mathbf{T}^n)\longrightarrow S/(\mathbf{T}^m)
  \qquad (0\leq n\leq m)
\]
a multiplicação por $(\mathrm{T}_0\cdots\mathrm{T}_r)^{m-n}$.
Para $n\geq\sup(p_i)_{0\leq i\leq r}$, denotemos por $\xi_{\mathbf{p}}^{(n)}=\xi_{p_0,\dots,p_r}^{(n)}$ a classe de $\mathrm{T}_0^{n-p_0}\cdots\mathrm{T}_r^{n-p_r}$ módulo $(\mathbf{T}^n)$.
Então $\vphi_{nm}(\xi_{\mathbf{p}}^{(n)})=\xi_{\mathbf{p}}^{(m)}$, de modo que estes elementos têm a mesma imagem canônica $\xi_{\mathbf{p}}=\xi_{p_0,\dots,p_r}$ no limite indutivo $\HH^{r+1}((\mathbf{T}),S)$.
Pela definição do grau em (\hyperref[III.2.1.2]{2.1.2}), o grau de $\xi_{\mathbf{p}}$ é $-|\mathbf{p}|=-(p_0+p_1+\cdots+p_r)$.
É claro que os $\xi_{\mathbf{p}}^{(n)}$ com $0<p_i\leq n$ para $0\leq i\leq r$ formam uma base de $S/(\mathbf{T}^n)$.
Deduz-se, então, imediatamente de \sref{III.2.1.8}:
\end{env}

\begin{corollary}[2.1.10]
\label{III.2.1.10}
Com a notação de \sref{III.2.1.8},
\[
\label{III.2.1.10.1}
  \HH^i((\mathbf{T}),S)=0
  \qquad\text{para }i\neq r+1,
  \tag{2.1.10.1}
\]
e $\HH^{r+1}((\mathbf{T}),S)$ é um $A$-módulo livre com base nos elementos $\xi_{p_0,\dots,p_r}$ tais que $p_i>0$ para $0\leq i\leq r$.
\end{corollary}

\begin{env}[2.1.11]
\label{III.2.1.11}
\emph{Observação.} ---
Seja $N$ um $A$-módulo qualquer e seja $M=S\otimes_A N$.
O raciocínio de \sref{III.2.1.8} mostra mais geralmente que
\[
\label{III.2.1.11.1}
  \HH^i(\mathbf{T}^n,M)=0
  \qquad\text{se }i\neq r+1,
  \tag{2.1.11.1}
\]
\[
\label{III.2.1.11.2}
  \HH^{r+1}(\mathbf{T}^n,M)=(S/(\mathbf{T}^n))\otimes_A N.
  \tag{2.1.11.2}
\]
Com efeito, esta última fórmula deduz-se diretamente de (\hyperref[III.1.1.3.5]{1.1.3.5}).
Além disso, é claro que
\[
  M/(\mathrm{T}_0^nM+\cdots+\mathrm{T}_{i-1}^nM)
  =\bigl(S/(\mathrm{T}_0^nS+\cdots+\mathrm{T}_{i-1}^nS)\bigr)\otimes_A N,
\]
pois o ideal $\mathrm{T}_0^nS+\cdots+\mathrm{T}_{i-1}^nS$ é um somando direto do $A$-módulo $S$.
Isso permite aplicar \sref{III.1.1.4} a $M$ e dá (\hyperref[III.2.1.11.1]{2.1.11.1}).
\end{env}

Combinando \sref{III.2.1.10} e \sref{III.2.1.5}, obtém-se:

\begin{proposition}[2.1.12]
\label{III.2.1.12}
Seja $A$ um anel, seja $r$ um inteiro $>0$ e seja $X=\bb{P}_A^r$ \sref[II]{II.4.1.1}.
Então:
\begin{enumerate}
  \item[{\rm(i) }] $\HH^i(X,\sh{O}_X(*))=0$ para $i\neq0,r$.
  \item[{\rm(ii) }] O homomorfismo canônico $\alpha:S\to\HH^0(X,\sh{O}_X(*))$ \sref[II]{II.2.6.2} é bijectivo.
\oldpage[III]{99}
  \item[{\rm(iii) }] $\HH^r(X,\sh{O}_X(*))$ é um $A$-módulo livre com base nos elementos $\xi_{p_0,\dots,p_r}$, onde $p_i>0$ para $0\leq i\leq r$; o elemento $\xi_{p_0,\dots,p_r}$ tem grau $-|\mathbf{p}|=-(p_0+\cdots+p_r)$, e
  \[
    \mathrm{T}_i\xi_{p_0,\dots,p_r}
    =\xi_{p_0,\dots,p_{i-1},p_i-1,p_{i+1},\dots,p_r}.
  \]
\end{enumerate}
\end{proposition}

\begin{proof}
Na sequência exata (\hyperref[III.2.1.5.2]{2.1.5.2}) aplicada a $M=S=A[\mathrm{T}_0,\dots,\mathrm{T}_r]$, temos $\HH^0((\mathbf{T}),S)=0$ e $\HH^1((\mathbf{T}),S)=0$ por \sref{III.2.1.10.1}; além disso, a Proposição~\sref{III.2.1.5} é aplicada com $U=X$, já que $X$ é a união dos $\mathrm{D}_+(\mathrm{T}_i)$ \sref[II]{II.2.3.14}.
Resta identificar a aplicação $S\to\HH^0(X,\sh{O}_X(*))$ de (\hyperref[III.2.1.5.2]{2.1.5.2}) com a aplicação canônica $\alpha$.
Isso se deduz da identificação canônica de $\HH^0(U,\sh{O}_X(*))$ e $\HH^0(\mathfrak{U},\sh{O}_X(*))$.
\end{proof}

\begin{corollary}[2.1.13]
\label{III.2.1.13}
Os únicos pares $(i,n)$ para os quais pode ocorrer que $\HH^i(X,\sh{O}_X(n))\neq0$ são $i=0$, $n\geq0$, e $i=r$, $n\leq-(r+1)$.
\end{corollary}

Se $A\neq0$, então temos efetivamente $\HH^i(X,\sh{O}_X(n))\neq0$ para todos os pares listados em \sref{III.2.1.13}.
Isto deduz-se de \sref{III.2.1.12}, já que $S_n\neq0$ para todo grau $n\geq0$.

Nas aplicações deste capítulo, utilizaremos principalmente o seguinte resultado menos preciso.

\begin{corollary}[2.1.14]
\label{III.2.1.14}
Os $A$-módulos $\HH^i(X,\sh{O}_X(n))$ são livres de tipo finito; se $i>0$, anulam-se para $n>0$.
\end{corollary}

\begin{proposition}[2.1.15]
\label{III.2.1.15}
Seja $Y$ um pré-esquema, seja $\sh{E}$ um $\sh{O}_Y$-módulo localmente livre de posto $r+1$, seja $X=\bb{P}(\sh{E})$ o fibrado projetivo definido por $\sh{E}$, e seja $f:X\to Y$ o morfismo estrutural.
Os únicos valores de $i$ e $n$ para os quais pode ocorrer que $\RR^if_*(\sh{O}_X(n))\neq0$ são $i=0$, $n\geq0$, e $i=r$, $n\leq-(r+1)$.
Além disso, o homomorfismo canônico \sref[II]{II.3.3.2}
\[
  \alpha:\bb{S}_{\sh{O}_Y}(\sh{E})
  \longrightarrow \Gamma_*(\sh{O}_X)
  =\RR^0f_*(\sh{O}_X(*))
  =\bigoplus_{n\in\bb{Z}}f_*(\sh{O}_X(n))
\]
é um isomorfismo.
\end{proposition}

\begin{proof}
A questão é local sobre $Y$, portanto podemos supor que $Y$ é afim, de anel $A$, e que $\sh{E}=\widetilde{E}$, onde $E=A^{r+1}$.
Então, ficamos imediatamente reduzidos a \sref{III.2.1.12}, tendo em conta \sref{III.1.4.11}.
\end{proof}

\begin{env}[2.1.16]
\label{III.2.1.16}
\emph{Observação.} ---
Mais adiante, completaremos os resultados de \sref{III.2.1.15} demonstrando as seguintes proposições.
Ponhamos
\[
  \boldsymbol{\omega}=f^*\bigl(\bigwedge^{r+1}\sh{E}\bigr)(-r-1),
\]
que é um $\sh{O}_X$-módulo invertível.
Então:
\begin{enumerate}
  \item[{\rm(i) }] Existe um isomorfismo canônico
  \[
  \label{III.2.1.16.1}
    \rho:\sh{O}_Y \simto \RR^rf_*(\boldsymbol{\omega}).
    \tag{2.1.16.1}
  \]
  \item[{\rm(ii) }] A acoplamento dado pelo produto cup \sref[0]{0.12.2.2}
  \[
  \label{III.2.1.16.2}
    \RR^rf_*(\sh{O}_X(n))
    \times \RR^0f_*(\boldsymbol{\omega}(-n))
    \longrightarrow \RR^rf_*(\boldsymbol{\omega}),
    \tag{2.1.16.2}
  \]
  composto por $\rho^{-1}$, define um \emph{isomorfismo} de $\RR^rf_*(\sh{O}_X(n))$ sobre o \emph{dual} do $\sh{O}_Y$-módulo localmente livre
  \[
    \RR^0f_*(\boldsymbol{\omega}(-n))
    =\bigl(\bigwedge^{r+1}\sh{E}\bigr)
      \otimes_{\sh{O}_Y}
      \bigl(\bb{S}_{\sh{O}_Y}(\sh{E})\bigr)_{-n}.
  \]
\end{enumerate}
\end{env}

\subsection{O teorema fundamental dos morfismos projetivos}
\label{subsection:III.2.2}
\oldpage[III]{100}

\begin{theorem}[2.2.1]
\label{III.2.2.1}
(Serre).
Seja $Y$ um pré-esquema noetheriano, seja $f:X\to Y$ um morfismo próprio e seja $\sh{L}$ um $\sh{O}_X$-módulo invertível amplo em relação a $f$.
Para todo $\sh{O}_X$-módulo $\sh{F}$, ponhamos
\[
  \sh{F}(n)=\sh{F}\otimes_{\sh{O}_X}\sh{L}^{\otimes n}
  \qquad(n\in\bb{Z}).
\]
Então, para todo $\sh{O}_X$-módulo coerente $\sh{F}$:
\begin{enumerate}
  \item[{\rm(i) }] os $\sh{O}_Y$-módulos $\RR^qf_*(\sh{F})$ são coerentes;
  \item[{\rm(ii) }] existe um inteiro $N$ tal que, para $n\geq N$,
  \[
    \RR^qf_*(\sh{F}(n))=0
    \qquad\text{Para tudo }q>0;
  \]
  \item[{\rm(iii) }] existe um inteiro $N$ tal que, para $n\geq N$, o homomorfismo canônico
  \[
    f^*\bigl(f_*(\sh{F}(n))\bigr)\longrightarrow\sh{F}(n)
  \]
  é sobrejetivo.
\end{enumerate}
\end{theorem}

\begin{proof}
Observemos primeiro que, se o teorema for verdadeiro depois de substituir $\sh{L}$ por $\sh{L}^{\otimes d}$, onde $d>0$, então será verdadeiro na sua forma original.
Com efeito, escrevendo
\[
  \sh{F}(n)
  =\bigl(\sh{F}\otimes\sh{L}^{\otimes r}\bigr)
   \otimes\sh{L}^{\otimes hd},
   \qquad h>0,\quad 0\leq r<d,
\]
A hipótese dá, para cada $r$, um inteiro $N_r$ tal que (ii) e (iii) valem para $\sh{F}\otimes\sh{L}^{\otimes r}$ sempre que $h\geq N_r$.
Tomando para $N$ o maior dos inteiros $dN_r$, as afirmações (ii) e (iii) valem, portanto, para $n\geq N$.

Podemos, portanto, supor que $\sh{L}$ é muito amplo relativamente a $f$ \sref[II]{II.4.6.11}.
Existe então uma imersão aberta dominante sobre $Y$
\[
  i:X\longrightarrow P,
  \qquad P=\Proj(\sh{S}),
\]
onde $\sh{S}$ é uma $\sh{O}_Y$-álgebra quase-coerente graduada em graus positivos, $\sh{S}_1$ é de tipo finito e gera $\sh{S}$, e $\sh{L}$ é isomorfo a $i^*(\sh{O}_P(1))$ \sref[II]{II.4.4.7}.
Como $f$ é próprio, $i$ também é próprio \sref[II]{II.5.4.4} e, portanto, é um isomorfismo $X\simto P$.
Assim, ficamos reduzidos ao caso $X=\Proj(\sh{S})$ e $\sh{L}=\sh{O}_X(1)$.
O teorema é então consequência da seguinte proposição.
\end{proof}

\begin{proposition}[2.2.2]
\label{III.2.2.2}
Seja $A$ um anel noetheriano e seja $S$ uma $A$-álgebra graduada em graus positivos tal que o $A$-módulo $S_1$ possui um sistema de $r+1$ geradores e gera a álgebra $S$.
Ponhamos $X=\Proj(S)$.
Para todo $\sh{O}_X$-módulo coerente $\sh{F}$:
\begin{enumerate}
  \item[{\rm(i) }] os $A$-módulos $\HH^q(X,\sh{F})$ são do tipo finito;
  \item[{\rm(ii) }] $\HH^q(X,\sh{F})=0$ para $q>r$;
  \item[{\rm(iii) }] existe um inteiro $N$ tal que, para $n\geq N$,
  \[
    \HH^q(X,\sh{F}(n))=0
    \qquad\text{para todo }q>0;
  \]
  \item[{\rm(iv) }] existe um inteiro $N$ tal que, para $n\geq N$, $\sh{F}(n)$ é gerado por suas seções sobre $X$.
\end{enumerate}
\end{proposition}

\begin{proof}
Primeiro vamos mostrar como \sref{III.2.2.2} implica \sref{III.2.2.1}.
No caso particular $X=\Proj(\sh{S})$ reduzido anteriormente \sref{III.2.2.1}, $Y$ é quase compacto.
Pode, portanto, ser recoberto por um número finito de abertos afins $U_\alpha$, de anéis noetherianos, tais que a restrição de $\sh{S}_1$ a cada $U_\alpha$ seja gerada por um número finito de seções sobre $U_\alpha$.
Uma vez conhecida \sref{III.2.2.2}, basta tomar, nas partes (ii) e (iii) de \sref{III.2.2.1}, o maior dos inteiros correspondentes aos $U_\alpha$, tendo em conta \sref{III.1.4.11} e \sref[II]{II.3.4.7}.

Para provar \sref{III.2.2.2}, observemos que $X$ se identifica com um subesquema fechado de $P=\bb{P}_A^r$ \sref[II]{II.3.6.2}.
Além disso, se $j:X\to P$ é a injeção canônica, então $j_*(\sh{F})$ é um $\sh{O}_P$-módulo coerente e
\[
  j_*(\sh{F}(n))=(j_*\sh{F})(n)
\]
por \sref[II]{II.3.4.5} e \sref[II]{II.3.5.2}.
Por força de (G, II, Corolário do Teorema~4.9.1), ficamos, portanto, reduzidos a provar \sref{III.2.2.2} no caso particular
\[
  X=\bb{P}_A^r,
  \qquad S=A[T_0,\ldots,T_r].
\]

\oldpage[III]{101}

Como $X$ é coberto pelos $r+1$ abertos afins $D_+(T_i)$, a afirmação (ii) deduz-se de \sref{III.1.4.12}.
A afirmação (iv) já foi demonstrada em \sref[II]{II.2.7.9}.

Agora vamos demonstrar simultaneamente (i) e (iii).
Eles valem para $\sh{F}=\sh{O}_X(m)$ por \sref{III.2.1.13} e, portanto, também quando $\sh{F}$ é a soma direta de um número finito de $\sh{O}_X$-módulos da forma $\sh{O}_X(m_j)$.
São trivialmente verdadeiras para $q>r$ por (ii).
Procedemos por indução descendente sobre $q$.

O feixe $\sh{F}$ é isomorfo a um quociente de uma soma direta $\sh{E}$ de um número finito de feixes $\sh{O}_X(m_j)$ \sref[II]{II.2.7.10}.
Existe, portanto, uma sequência exata
\[
  0\longrightarrow\sh{R}\longrightarrow\sh{E}
   \longrightarrow\sh{F}\longrightarrow0,
\]
onde $\sh{R}$ é coerente \sref[0]{0.5.3.3} e $\sh{E}$ satisfaz (i) e (iii).
Como a mudança de grau é exata, para todo $n\in\bb{Z}$ temos também uma sequência exata
\[
  0\longrightarrow\sh{R}(n)\longrightarrow\sh{E}(n)
   \longrightarrow\sh{F}(n)\longrightarrow0,
\]
e, portanto, uma sequência exata de cohomologia
\[
  \HH^{q-1}(X,\sh{E}(n))
  \longrightarrow\HH^{q-1}(X,\sh{F}(n))
  \longrightarrow\HH^q(X,\sh{R}(n)).
\]
Como $\sh{E}(n)$ é uma soma direta de feixes $\sh{O}_X(n+m_j)$ \sref[II]{II.2.5.14}, o módulo $\HH^{q-1}(X,\sh{E}(n))$ é do tipo finito.
Analogamente, o módulo $\HH^q(X,\sh{R}(n))$ é do tipo finito pela hipótese de indução.
Como $A$ é noetheriano, deduz-se que $\HH^{q-1}(X,\sh{F}(n))$ é do tipo finito, o que demonstra (i).

Pela hipótese de indução podemos escolher $N$ de modo que $\HH^q(X,\sh{R}(n))=0$ para $n\geq N$.
Podemos também escolher $N$ para que $\HH^{q-1}(X,\sh{E}(n))=0$ para $n\geq N$, já que $\sh{E}$ satisfaz (iii).
A sequência exata de cohomologia dá então
\[
  \HH^{q-1}(X,\sh{F}(n))=0
  \qquad(n\geq N),
\]
o que conclui a demonstração.
\end{proof}

\begin{corollary}[2.2.3]
\label{III.2.2.3}
Sob as hipóteses de \sref{III.2.2.1}, seja
\[
  \sh{F}\longrightarrow\sh{G}\longrightarrow\sh{H}
\]
uma sequência exata de $\sh{O}_X$-módulos coerentes.
Existe um inteiro $N$ tal que, para $n\geq N$, a sequência
\[
  f_*(\sh{F}(n))
  \longrightarrow f_*(\sh{G}(n))
  \longrightarrow f_*(\sh{H}(n))
\]
é exata.
\end{corollary}

\begin{proof}
Sejam $\sh{F}'$, $\sh{G}'$ e $\sh{G}''$, respectivamente, o núcleo, a imagem e o conúcleo de $\sh{F}\to\sh{G}$.
Então $\sh{G}'$ é o núcleo e $\sh{G}''$ a imagem de $\sh{G}\to\sh{H}$; seja $\sh{H}''$ o conúcleo deste último homomorfismo.
Todos estes $\sh{O}_X$-módulos são coerentes \sref[0]{0.5.3.4}.
Como a mudança de grau é exata, basta provar que, para $n$ suficientemente grande, cada uma das sequências
\begin{align*}
  0&\longrightarrow f_*(\sh{F}'(n))
    \longrightarrow f_*(\sh{F}(n))
    \longrightarrow f_*(\sh{G}'(n))
    \longrightarrow0,\\
  0&\longrightarrow f_*(\sh{G}'(n))
    \longrightarrow f_*(\sh{G}(n))
    \longrightarrow f_*(\sh{G}''(n))
    \longrightarrow0,\\
  0&\longrightarrow f_*(\sh{G}''(n))
    \longrightarrow f_*(\sh{H}(n))
    \longrightarrow f_*(\sh{H}''(n))
    \longrightarrow0
\end{align*}
é exata.
Consequentemente, podemos supor que
\[
  0\longrightarrow\sh{F}\longrightarrow\sh{G}
   \longrightarrow\sh{H}\longrightarrow0
\]
é exata.
A sequência de cohomologia é então
\[
  0\longrightarrow f_*(\sh{F}(n))
  \longrightarrow f_*(\sh{G}(n))
  \longrightarrow f_*(\sh{H}(n))
  \longrightarrow \RR^1f_*(\sh{F}(n))
  \longrightarrow\cdots,
\]
e a conclusão deduz-se de \sref{III.2.2.1}, (ii).
\end{proof}

\oldpage[III]{102}

\begin{corollary}[2.2.4]
\label{III.2.2.4}
Seja $Y$ um pré-esquema noetheriano, seja $f:X\to Y$ um morfismo de tipo finito e seja $\sh{L}$ um $\sh{O}_X$-módulo invertível amplo em relação a $f$.
Para todo $\sh{O}_X$-módulo $\sh{F}$, ponhamos
\[
  \sh{F}(n)=\sh{F}\otimes_{\sh{O}_X}\sh{L}^{\otimes n}
  \qquad(n\in\bb{Z}).
\]
Seja
\[
  \sh{F}\longrightarrow\sh{G}\longrightarrow\sh{H}
\]
uma sequência exata de $\sh{O}_X$-módulos coerentes tal que os suportes de $\sh{F}$ e $\sh{H}$ sejam próprios sobre $Y$ \sref[II]{II.5.4.10}.
Então há um inteiro $N$ tal que, para $n\geq N$, a sequência
\[
  f_*(\sh{F}(n))
  \longrightarrow f_*(\sh{G}(n))
  \longrightarrow f_*(\sh{H}(n))
\]
é exata.
\end{corollary}

\begin{proof}
O raciocínio do início de \sref{III.2.2.1} mostra que, se o corolário for verdadeiro para $\sh{L}^{\otimes d}$, onde $d>0$, então será verdadeiro para $\sh{L}$.
Podemos, portanto, supor que $\sh{L}$ é muito amplo em relação a $f$ \sref[II]{II.4.6.11}.
Podemos então identificar $X$ com um subconjunto aberto de um $Y$-pré-esquema
\[
  Z=\Proj(\sh{S}),
\]
onde $\sh{S}$ é uma $\sh{O}_Y$-álgebra quase-coerente graduada em graus positivos, $\sh{S}_1$ é de tipo finito e gera $\sh{S}$, e
\[
  \sh{L}=i^*(\sh{O}_Z(1)),
\]
sendo $i:X\to Z$ a imersão canônica \sref[II]{II.4.4.7}.

Agora $\Supp(\sh{G})$ é fechado em $X$ e está contido em
\[
  \Supp(\sh{F})\cup\Supp(\sh{H});
\]
portanto, é próprio sobre $Y$.
Consequentemente, os suportes de $\sh{F}$, $\sh{G}$ e $\sh{H}$ são fechados em $Z$ \sref[II]{II.5.4.10}.
Assim,
\[
  \sh{F}'=i_*(\sh{F}),\qquad
  \sh{G}'=i_*(\sh{G}),\qquad
  \sh{H}'=i_*(\sh{H})
\]
são $\sh{O}_Z$-módulos coerentes, e a sequência
\[
  \sh{F}'\longrightarrow\sh{G}'\longrightarrow\sh{H}'
\]
é exata.
Se $g:Z\to Y$ é o morfismo estrutural, então $f=g\circ i$, e claramente
\[
  \sh{F}'(n)=i_*(\sh{F}(n)),
\]
com identidades análogas para $\sh{G}'$ e $\sh{H}'$.
A conclusão deduz-se aplicando \sref{III.2.2.3} a $\sh{F}'$, $\sh{G}'$ e $\sh{H}'$.
\end{proof}

\begin{env}[2.2.5]
\label{III.2.2.5}
\emph{Observações.}

(i) A afirmação (i) de \sref{III.2.2.1} continua sendo verdadeira se se supõe apenas que $Y$ é \emph{localmente noetheriano}.
Com efeito, a propriedade é manifestamente local sobre $Y$.
Além disso, as hipóteses de \sref{III.2.2.1} implicam que, para todo subconjunto aberto $U\subset Y$, a restrição de $f$ a $f^{-1}(U)$ é um morfismo projetivo
\[
  f^{-1}(U)\longrightarrow U
\]
\sref[II]{II.5.5.5}[iii], e $\sh{L}|_{f^{-1}(U)}$ é amplo para este morfismo \sref[II]{II.4.6.4}.

(ii) A afirmação (iii) de \sref{III.2.2.1} continua válida, como vimos, se se supõe apenas que $X$ é um esquema quase compacto, ou um pré-esquema cujo espaço subjacente é noetheriano, e que $f:X\to Y$ é quase compacto \sref[II]{II.4.6.8}.
É interessante observar, no entanto, que mesmo quando $Y$ é o \emph{espectro de um corpo} $K$ e $f$ é \emph{quase-projetivo}, a afirmação (ii) de \sref{III.2.2.1} pode deixar de se cumprir.

Por exemplo, seja
\[
  X'=\Spec(K[T_0,\ldots,T_r])
\]
e seja $X$ a união dos subconjuntos abertos afins $D(T_i)$ de $X'$, para $0\leq i\leq r$.
Como a imersão $X\to X'$ é quase compacta, o morfismo estrutural $f:X\to Y$ é quase afim \sref[II]{II.5.1.10} e, portanto, $\sh{O}_X$ é muito amplo em relação a $f$ \sref[II]{II.5.1.6}.
Mas
\[
  \Gamma(X,\sh{O}_X)
  =\bigcap_{i=0}^r(K[T_0,\ldots,T_r])_{T_i}
  =K[T_0,\ldots,T_r]
\]
por \sref[I]{I.8.2.1.1}.
Deduz-se de \sref{III.1.4.3.1} e \sref{III.1.1.3.5} que
\[
  \HH^r(X,\sh{O}_X^{\otimes n})
  =\HH^r(X,\sh{O}_X)
  =A\neq0
\]
para todo $n$.
\end{env}


\subsection{Aplicação a feixes graduados de álgebras e de módulos}
\label{subsection:III.2.3}

\begin{theorem}[2.3.1]
\label{III.2.3.1}
Seja $Y$ um pré-esquema noetheriano, seja $\sh{S}$ uma $\sh{O}_Y$-álgebra quase-coerente de tipo finito, graduada em graus positivos, seja
\[
  X=\Proj(\sh{S}),
\]
e seja $q:X\to Y$ o morfismo estrutural.
Seja $\sh{M}$ um $\sh{S}$-módulo graduado quase-coerente que satisfaça a condição \textbf{(TF) }.
Existe então um inteiro $N$ tal que, para $n\geq N$, o homomorfismo canônico \sref[II]{II.8.14.5.1}
\[
  \alpha_n:\sh{M}_n
  \longrightarrow q_*\bigl(\shProj_0(\sh{M}(n))\bigr)
  =q_*\bigl((\shProj(\sh{M}))_n\bigr)
\]
\oldpage[III]{103}
é bijetivo.
Em outras palavras, o homomorfismo canônico
\[
  \alpha:\sh{M}\longrightarrow\bbGamma_*(\shProj(\sh{M}))
\]
é um \textbf{(TN) }-isomorfismo.
\end{theorem}

\begin{proof}
Podemos restringir-nos ao caso em que $\sh{M}$ é um $\sh{S}$-módulo de tipo finito \sref[II]{II.3.4.2}.
Como $Y$ é quase compacto, existe um inteiro $d>0$ tal que $\sh{S}^{(d)}$ é gerada pelo $\sh{O}_Y$-módulo quase-coerente $\sh{S}_d$, que é do tipo finito \sref[II]{II.3.1.10} e, portanto, coerente porque $Y$ é noetheriano.
Agora, $\sh{M}$ é a soma direta dos $\sh{M}^{(d,k)}$, para $0\leq k<d$, e cada $\sh{M}^{(d,k)}$ é um $\sh{S}^{(d)}$-módulo quase-coerente de tipo finito, por \sref[II]{II.2.1.6}[iii].
Como a questão é local sobre $Y$, basta manifestamente demonstrar que cada homomorfismo canônico
\[
  \alpha:\sh{M}^{(d,k)}
  \longrightarrow
  \bbGamma_*\bigl((\shProj(\sh{M}))^{(d,k)}\bigr)
\]
é um \textbf{(TN) }-isomorfismo.
Tendo em conta \sref[II]{II.8.14.13}, e em particular o diagrama \sref[II]{II.8.14.13.4}, ficamos assim reduzidos a demonstrar o teorema quando $\sh{S}$ é gerada por $\sh{S}_1$ e $\sh{S}_1$ é um $\sh{O}_Y$-módulo coerente.

Como $Y$ é noetheriano, o raciocínio usado no início da demonstração de \sref{III.2.2.2} mostra que podemos reduzir-nos ao caso
\[
  Y=\Spec(A),\qquad
  \sh{S}=\widetilde{S},\qquad
  \sh{M}=\widetilde{M},
\]
onde $A$ é um anel noetheriano, $S_1$ é um $A$-módulo de tipo finito e $M$ é um $S$-módulo graduado de tipo finito.
Vamos mostrar que basta então provar o teorema para $M=S$.
Com efeito, no caso geral existe uma sequência exata
\[
  L'\longrightarrow L\longrightarrow M\longrightarrow0,
\]
onde $L$ e $L'$ são somas diretas de módulos graduados da forma $S(m)$.
Se o resultado vale para $M=S$, também vale para $M=S(m)$ e, portanto, para $L$ e $L'$.
Consideremos o diagrama comutativo
\[
  \xymatrix{
    \widetilde{L'_n}\ar[r]\ar[d]_{\alpha_n} &
    \widetilde{L_n}\ar[r]\ar[d]^{\alpha_n} &
    \widetilde{M_n}\ar[r]\ar[d]^{\alpha_n} &
    0\\
    q_*(\widetilde{L'}(n))\ar[r] &
    q_*(\widetilde{L}(n))\ar[r] &
    q_*(\widetilde{M}(n))\ar[r] &
    0 .
  }
\]
\oldpage[III]{104}
A segunda fila é exata por \sref{III.2.2.3} quando $n$ é suficientemente grande.
O mesmo acontece com a primeira fila e, como as duas flechas verticais à esquerda são isomorfismos, também é a terceira.

Resta, então, provar o teorema para $M=S$.
Suponhamos que, em primeiro lugar,
\[
  S=A[T_0,\ldots,T_r],
\]
onde os $T_i$ são indeterminadas.
Neste caso, a afirmação é precisamente \sref{III.2.1.11}[ii].
No caso geral, $S$ identifica-se com um quociente de um anel
\[
  S'=A[T_0,\ldots,T_r]
\]
por um ideal graduado, de modo que $X$ é um subesquema fechado de
\[
  X'=\bb{P}_A^r
\]
\sref[II]{II.2.9.2}.
Se $j:X\to X'$ é a injeção canônica, então $j_*(\widetilde{S}(n))$ é o $\sh{O}_{X'}$-módulo
\[
  (\shProj(\widetilde{S}))(n),
\]
onde $S$ é considerado como $S'$-módulo graduado; isso se deduz imediatamente de \sref[II]{II.2.8.7}.
Como $j_*(\widetilde{S}(n))$ é um $\sh{O}_{X'}$-módulo que satisfaz \textbf{(TF) }, o homomorfismo canônico
\[
  \alpha_n:S_n\longrightarrow
  \Gamma\bigl(X',j_*(\widetilde{S}(n))\bigr)
\]
é bijetivo para todo $n$ suficientemente grande, pelo caso polinomial que acabamos de demonstrar.
Isso conclui a demonstração, já que
\[
  \Gamma\bigl(X',j_*(\widetilde{S}(n))\bigr)
  =\Gamma(X,\widetilde{S}(n)).
\]
\end{proof}

\begin{corollary}[2.3.2]
\label{III.2.3.2}
Sob as hipóteses de \sref{III.2.3.1}, seja
\[
  \sh{S}_X=\bigoplus_{n\in\bb{Z}}\sh{O}_X(n),
\]
e seja $\sh{F}$ um $\sh{S}_X$-módulo graduado quase-coerente de tipo finito.
Então $\bbGamma_*(\sh{F})$ satisfaz a condição \textbf{(TF) }.
\end{corollary}

\begin{proof}
Na demonstração de \sref{III.2.3.1} vimos que $X$, que é isomorfo a $\Proj(\sh{S}^{(d)})$ \sref[II]{II.3.1.8}, é de tipo finito sobre $Y$ \sref[II]{II.3.4.1}.
Deduz-se de \sref[II]{II.8.14.9} que $\sh{F}$ é isomorfo a um $\sh{S}_X$-módulo graduado da forma $\shProj(\sh{M})$, onde $\sh{M}$ é um $\sh{S}$-módulo graduado quase-coerente do tipo finito.
Por \sref{III.2.3.1}, $\bbGamma_*(\sh{F})$ é \textbf{(TN) }-isomorfo a $\sh{M}$ e, portanto, satisfaz \textbf{(TF) }.
\end{proof}

\begin{env}[2.3.3]
\label{III.2.3.3}
\emph{Escolio.}
Seja $Y$ um pré-esquema noetheriano, seja $\sh{S}$ uma $\sh{O}_Y$-álgebra graduada que satisfaz as hipóteses de \sref{III.2.3.1}, e seja $X=\Proj(\sh{S})$.
Seja $\cat{K}_{\sh{S}}$ a categoria abeliana dos $\sh{S}$-módulos graduados quase-coerentes que satisfazem \textbf{(TF) }, e seja $\cat{K}'_{\sh{S}}$ a subcategoria de $\cat{K}_{\sh{S}}$ formada pelos $\sh{S}$-módulos que satisfazem \textbf{(TN) }.
Por último, seja $\cat{K}_X$ a categoria dos $\sh{S}_X$-módulos graduados quase-coerentes $\sh{F}$ de tipo finito; como $\sh{S}_X$ é periódica \sref[II]{II.8.14.4} e \sref[II]{II.8.14.12}, isto equivale a exigir que os $\sh{F}_i$ sejam $\sh{O}_X$-módulos coerentes.

Por \sref[II]{II.8.14.8}, \sref[II]{II.8.14.10} e \sref{III.2.3.2}, os funtores
\[
  \sh{M}\longmapsto\shProj(\sh{M}),
  \qquad
  \sh{F}\longmapsto\bbGamma_*(\sh{F})
\]
definem uma equivalência (T, I, 1.2) entre a categoria quociente
\[
  \cat{K}_{\sh{S}}/\cat{K}'_{\sh{S}}
\]
(T, I, 1.11) e a categoria $\cat{K}_X$.
Quando $\sh{S}$ é gerada por $\sh{S}_1$, a categoria $\cat{K}_X$ pode ser substituída pela categoria dos $\sh{O}_X$-módulos coerentes \sref[II]{II.8.14.12}.
\end{env}

\begin{proposition}[2.3.4]
\label{III.2.3.4}
Seja $Y$ um pré-esquema noetheriano.

(i) Seja $\sh{S}$ uma $\sh{O}_Y$-álgebra quase-coerente de tipo finito, graduada em graus positivos.
Ponhamos
\[
  X=\Proj(\sh{S}),\qquad
  \sh{S}_X=\shProj(\sh{S})
  =\bigoplus_{n\in\bb{Z}}\sh{O}_X(n).
\]
Então $\sh{S}_X$ é uma $\sh{O}_X$-álgebra graduada quase-coerente periódica \sref[II]{II.8.14.12}, suas componentes homogêneas
\[
  (\sh{S}_X)_n=\sh{O}_X(n)
\]
são $\sh{O}_X$-módulos coerentes e, se $d>0$ é um período de $\sh{S}_X$, então
\[
  (\sh{S}_X)_d=\sh{O}_X(d)
\]
é $Y$-amplo como $\sh{O}_X$-módulo invertível.
Além disso, o homomorfismo canônico
\[
  \alpha:\sh{S}\longrightarrow\bbGamma_*(\sh{S}_X)
\]
é um \textbf{(TN) }-isomorfismo.

(ii) Reciprocamente, seja $q:X\to Y$ um morfismo projetivo e seja $\sh{S}'$ uma $\sh{O}_X$-álgebra graduada cujas componentes homogêneas $\sh{S}'_n$, para $n\in\bb{Z}$, são $\sh{O}_X$-módulos coerentes.
Suponhamos que $\sh{S}'$ tenha um período $d>0$ tal que $\sh{S}'_d$ seja um $\sh{O}_X$-módulo invertível amplo para $q$.
Então
\[
  \sh{S}=\bigoplus_{n\geq0}q_*(\sh{S}'_n)
\]
é uma $\sh{O}_Y$-álgebra quase-coerente de tipo finito, graduada em graus positivos, e existe um $Y$-isomorfismo
\[
  r:X\simto\Proj(\sh{S})
\]
tal que $r^*(\shProj(\sh{S}))$ é isomorfo a $\sh{S}'$ como $\sh{O}_X$-álgebra graduada.
\end{proposition}

\begin{proof}
(i) Quase todas as afirmações já foram demonstradas; a última é um caso particular de \sref{III.2.3.2}.
A periodicidade de $\sh{S}_X$ foi demonstrada em \sref[II]{II.8.14.14}, e a afirmação de que existe um período $d>0$ para o qual $\sh{O}_X(d)$ é invertível e $Y$-amplo é precisamente \sref[II]{II.4.6.18}.
Finalmente, para $0\leq k<d$, $(\sh{S}_X)^{(d,k)}$ é um $(\sh{S}_X)^{(d)}$-módulo de tipo finito \sref[II]{II.8.14.14}.
Deduz-se que cada $(\sh{S}_X)_n$ é um $\sh{O}_X$-módulo quase-coerente de tipo finito \sref[II]{II.2.1.6}[ii]; como $\sh{O}_X$ é coerente, cada $(\sh{S}_X)_n$ também o é.

(ii) Substituindo o período $d$ por um de seus múltiplos, podemos supor que
\[
  \sh{L}=\sh{S}'_d
\]
é muito amplo em relação a $q$ \sref[II]{II.4.6.11}.
Por hipótese,
\[
  \sh{S}'^{(d)}
  =\bigoplus_{n\in\bb{Z}}\sh{L}^{\otimes n},
\]
e, portanto,
\[
  \sh{S}^{(d)}
  =\bigoplus_{n\geq0}q_*(\sh{L}^{\otimes n}).
\]
Por \sref[II]{II.3.1.8} e \sref[II]{II.3.2.9}, existe um $Y$-isomorfismo
\[
  s:X'=\Proj(\sh{S})\simto
  X''=\Proj(\sh{S}^{(d)})
\]
tal que
\oldpage[III]{105}
\[
  s^*(\sh{O}_{X''}(n))=\sh{O}_{X'}(nd)
\]
para todo $n$ \sref[II]{II.3.5.2}.
Assim, estabeleceremos a existência de um $Y$-isomorfismo $X\simto X'$ demonstrando a seguinte afirmação.

\begin{env}[2.3.4.1]
\label{III.2.3.4.1}
Seja $Y$ um pré-esquema noetheriano, seja $q:X\to Y$ um morfismo projetivo e seja $\sh{L}$ um $\sh{O}_X$-módulo invertível muito amplo para $q$.
Então
\[
  \sh{S}=\bigoplus_{n\geq0}q_*(\sh{L}^{\otimes n})
\]
é uma $\sh{O}_Y$-álgebra graduada quase-coerente de tipo finito tal que
\[
  \sh{S}_n=\sh{S}_1^n
\]
para todo $n$ suficientemente grande, e existe um $Y$-isomorfismo
\[
  r:X\simto P=\Proj(\sh{S})
\]
tal que
\[
  \sh{L}=r^*(\sh{O}_P(1)).
\]
\end{env}

Como $q$ é projetivo, deduz-se de \sref[II]{II.5.4.4} e \sref[II]{II.4.4.7} que existe um $Y$-isomorfismo
\[
  r':X\simto P'=\Proj(\sh{T}),
\]
onde $\sh{T}$ é uma $\sh{O}_Y$-álgebra quase-coerente, $\sh{T}_1$ é de tipo finito e gera $\sh{T}$, e
\[
  \sh{L}=r'^*(\sh{O}_{P'}(1)).
\]
Se $q':P'\to Y$ é o morfismo estrutural, então
\[
  \sh{S}=\bigoplus_{n\geq0}q'_*(\sh{O}_{P'}(n)).
\]
Por \sref{III.2.3.1}, o homomorfismo canônico
\[
  \alpha_n:\sh{T}_n\longrightarrow
  \sh{S}_n=q'_*(\sh{O}_{P'}(n))
\]
é bijetivo para todo $n$ suficientemente grande.
Como $\sh{T}_n=\sh{T}_1^n$, deduz-se \emph{a fortiori} que $\sh{S}_n=\sh{S}_1^n$ para todo $n$ suficientemente grande.

O homomorfismo canônico $\alpha:\sh{T}\to\sh{S}$ de $\sh{O}_Y$-álgebras graduadas é um \textbf{(TN) }-isomorfismo.
Consequentemente,
\phantomsection\label{III.2.3.5.1}
\[
  \Phi=\Proj(\alpha):\Proj(\sh{S})\longrightarrow\Proj(\sh{T})
\]
é um isomorfismo \sref[II]{II.3.6.1}.
Além disso, \sref[II]{II.3.5.2} e o fato de que os $\sh{T}$-módulos graduados obtidos de $\sh{S}(n)$ por restrição de escalares são \textbf{(TN) }-isomorfos a $\sh{T}(n)$ dão
\[
  \Phi_*(\sh{O}_P(n))=\sh{O}_{P'}(n)
\]
para todo $n$ \sref[II]{II.3.4.2}.
Só resta mostrar que $\sh{S}$ é de tipo finito sobre $\sh{O}_Y$.
Os módulos
\[
  \sh{S}_n=q'_*(\sh{O}_{P'}(n))
\]
são coerentes por \sref{III.2.2.1}; e, se $\sh{S}_n=\sh{S}_1^n$ para $n\geq n_0$, então $\sh{S}$ é gerada pelo módulo coerente
\[
  \bigoplus_{i\leq n_0}\sh{S}_i.
\]
Isto demonstra a afirmação \sref[I]{I.9.6.2}.

Voltemos à demonstração de \sref{III.2.3.4}, mantendo a notação.
Demonstramos que existe um $Y$-isomorfismo
\[
  r'':X\simto X''
\]
tal que
\[
  r''_*(\sh{L}^{\otimes n})=\sh{O}_{X''}(n)
\]
para todo $n\in\bb{Z}$.
Seja $q'':X''\to Y$ o morfismo estrutural.
A álgebra $\sh{S}'$ é a soma direta dos $\sh{S}'^{(d)}$-módulos graduados $\sh{S}'^{(d,k)}$.
Cada um deles é um $\sh{S}'^{(d)}$-módulo quase-coerente de tipo finito, pela periodicidade de $\sh{S}'$ e pela hipótese de que os $\sh{S}'_n$ são $\sh{O}_X$-módulos de tipo finito \sref[II]{II.8.14.12}.
Ponhamos
\[
  \sh{F}^{(k)}=r''_*(\sh{S}'^{(d,k)}),
\]
de modo que os $\sh{F}^{(k)}$ são $\sh{S}_{X''}$-módulos graduados quase-coerentes de tipo finito.
Por \sref[II]{II.8.14.8}, o homomorfismo canônico
\[
  \beta:\shProj\bigl(\bbGamma_*(\sh{F}^{(k)})\bigr)
  \longrightarrow\sh{F}^{(k)}
\]
é um isomorfismo de $\sh{S}_{X''}$-módulos.
Por outro lado,
\[
  q''_*\bigl((\sh{F}^{(k)})_n\bigr)
  =q_*\bigl((\sh{S}'^{(d,k)})_n\bigr),
\]
e para $n\geq0$ o último $\sh{O}_Y$-módulo é, por definição, igual a $(\sh{S}^{(d,k)})_n$.
Em outras palavras, a injeção canônica
\[
  \sh{S}^{(d,k)}
  \longrightarrow\bbGamma_*(\sh{F}^{(k)})
\]
é um \textbf{(TN) }-isomorfismo.
Portanto,
\[
  \shProj(\sh{S}^{(d,k)})
  =\shProj\bigl(\bbGamma_*(\sh{F}^{(k)})\bigr),
\]
e, consequentemente,
\[
  r''{}^*\bigl(\shProj(\sh{S}^{(d,k)})\bigr)
  \simeq\sh{S}'^{(d,k)}.
\]
Resta observar que, salvo isomorfismo canônico,
\[
  \shProj(\sh{S}^{(d,k)})
  =s_*\bigl((\shProj(\sh{S}))^{(d,k)}\bigr)
\]
\sref[II]{II.8.14.13.1}; isto demonstra que $r^*(\shProj(\sh{S}))$ é isomorfo a $\sh{S}'$.

Por fim, por \sref{III.2.3.2}, cada $\bbGamma_*(\sh{F}^{(k)})$ satisfaz a condição \textbf{(TF) } e, portanto, também a satisfaz cada $\sh{S}^{(d,k)}$.
Além disso, os $\sh{S}_n=q_*(\sh{S}'_n)$ são coerentes por \sref{III.2.2.1}, porque os $\sh{S}'_n$ são coerentes.
Deduz-se imediatamente que os $\sh{S}^{(d,k)}$ são $\sh{S}^{(d)}$-módulos de tipo finito.
Como \sref{III.2.3.4.1} mostra que $\sh{S}^{(d)}$ é uma $\sh{O}_Y$-álgebra de tipo finito, $\sh{S}$ é também uma $\sh{O}_Y$-álgebra de tipo finito.
\end{proof}

\oldpage[III]{106}
\begin{proposition}[2.3.5]
\label{III.2.3.5}
Seja $Y$ um pré-esquema noetheriano integral, seja $X$ um pré-esquema integral e seja
\[
  f:X\longrightarrow Y
\]
um morfismo projetivo birracional.
Existe então um ideal fracionário coerente
\[
  \sh{J}\subset\sh{R}(Y)
\]
\sref[II]{II.8.1.2} tal que $X$ é $Y$-isomorfo ao pré-esquema obtido por explosão de $\sh{J}$ \sref[II]{II.8.1.3}.
Além disso, existe um subconjunto aberto $U$ de $Y$ tal que a restrição de $f$ a $f^{-1}(U)$ é um isomorfismo de $f^{-1}(U)$ sobre $U$ (cf. \ \textbf{I}, 6.5.5), e tal que $\sh{J}|U$ é invertível.
\end{proposition}

\begin{proof}
Existe um $\sh{O}_X$-módulo invertível $\sh{L}$ muito amplo para $f$ \sref[II]{II.4.4.2} e \sref[II]{II.5.3.2}.
Aplicando \sref{III.2.3.4.1}, identificamos $X$ com $\Proj(\sh{S})$, onde
\[
  \sh{S}=\bigoplus_{n\geq0}f_*(\sh{L}^{\otimes n}).
\]
Os módulos $f_*(\sh{L}^{\otimes n})$ são sem torção sobre $\sh{O}_Y$ \sref[I]{I.7.4.5}.
Assim, $\sh{S}$ identifica-se canonicamente com um sub-$\sh{O}_Y$-módulo de
\[
  \sh{S}\otimes_{\sh{O}_Y}\sh{R}(Y)
\]
\sref[I]{I.7.4.1}.
O último é um feixe simples \sref[I]{I.7.3.6} e é, portanto, determinado por sua restrição a um subconjunto aberto não vazio.

Escolhamos um subconjunto aberto não vazio $U'\subset U$ sobre o qual $\sh{L}|f^{-1}(U')$ seja isomorfo a $\sh{O}_X|f^{-1}(U')$ e, portanto, a $\sh{O}_Y|U'$.
As restrições $f_*(\sh{L}^{\otimes n})|U'$ são então isomorfas a $\sh{O}_Y|U'$.
Segue-se que
\[
  \sh{S}\otimes_{\sh{O}_Y}\sh{R}(Y)
  \simeq\sh{R}(Y)[T],
\]
onde $T$ é uma indeterminada.
Por \sref{III.2.3.4.1}, $\sh{S}$ é \textbf{(TN) }-isomorfa à sub-$\sh{O}_Y$-álgebra gerada pela imagem canônica de $f_*(\sh{L})$ nesta álgebra de polinômios.
Sob a identificação mostrada, essa imagem é $\sh{J}\cdot T$, onde $\sh{J}$ é um sub-$\sh{O}_Y$-módulo coerente de $\sh{R}(Y)$ \sref{III.2.2.1}.
Sua restrição a $U'$ é isomorfa a $\sh{O}_Y|U'$ e, reduzindo $U$ se necessário, $\sh{J}|U$ é invertível.
Consequentemente,
\[
  \sh{S}\quad\text{é \textbf{(TN)}-isomorfa a}\quad
  \bigoplus_{n\geq0}\sh{J}^n,
\]
o que conclui a demonstração.
\end{proof}

\begin{corollary}[2.3.6]
\label{III.2.3.6}
Sob as hipóteses de \sref{III.2.3.5}, suponhamos ainda que, para todo sub-$\sh{O}_Y$-módulo coerente não nulo,
\[
  \sh{J}\subset\sh{R}(Y)
\]
existe um $\sh{O}_Y$-módulo invertível $\sh{L}$ tal que
\[
  \Gamma\bigl(Y,\sh{L}\otimes_{\sh{O}_Y}
  \shHom(\sh{J},\sh{O}_Y)\bigr)\neq0.
\]
Então, na declaração de \sref{III.2.3.5}, pode-se tomar $\sh{J}$ como um ideal de $\sh{O}_Y$.
Esta condição adicional é sempre satisfeita se houver um $\sh{O}_Y$-módulo amplo.
\end{corollary}

\begin{proof}
De fato, por \sref[0]{0.5.4.2},
\[
  \sh{L}\otimes\shHom(\sh{J},\sh{O}_Y)
  =\shHom\bigl(\sh{L}^{-1},\shHom(\sh{J},\sh{O}_Y)\bigr)
  =\shHom(\sh{J}\otimes\sh{L}^{-1},\sh{O}_Y).
\]
A hipótese diz, portanto, que existe um homomorfismo não nulo
\[
  u:\sh{J}\otimes\sh{L}^{-1}\longrightarrow\sh{O}_Y.
\]
Para todo $y\in Y$, o módulo $(\sh{J}\otimes\sh{L}^{-1})_y$ identifica-se com um sub-$\sh{O}_y$-módulo do corpo de frações $\sh{R}(Y)_y$ de $\sh{O}_y$ \sref[I]{I.7.1.5}.
Assim, $u_y$ é necessariamente injetivo, e $u$ identifica $\sh{J}\otimes\sh{L}^{-1}$ com um ideal $\sh{J}'$ de $\sh{O}_Y$.
Mas
\[
  \Proj\left(\bigoplus_{n\geq0}\sh{J}^n\right)
  \quad\text{e}\quad
  \Proj\left(\bigoplus_{n\geq0}
  (\sh{J}\otimes\sh{L}^{-1})^n\right)
\]
são $Y$-isomorfos \sref[II]{II.3.1.8}; isto demonstra a primeira afirmação.

Para a segunda, observe-se que
\[
  \sh{F}=\shHom_{\sh{O}_Y}(\sh{J},\sh{O}_Y)
\]
é coerente e não nulo, pois existe um subconjunto aberto $U$ de $Y$ sobre o qual $\sh{J}|U$ é invertível.
Se $\sh{L}$ é um $\sh{O}_Y$-módulo amplo, existe um inteiro $n$ tal que
\[
  \sh{F}(n)=\sh{F}\otimes\sh{L}^{\otimes n}
\]
é gerado por suas seções sobre $Y$ \sref[II]{II.4.5.5}.
\emph{A fortiori},
\[
  \Gamma(Y,\sh{F}(n))\neq0,
\]
o que demonstra a afirmação.
\end{proof}

\begin{corollary}[2.3.7]
\label{III.2.3.7}
Sejam $X$ e $Y$ esquemas integrais projetivos sobre um corpo $k$, e seja
\[
  f:X\longrightarrow Y
\]
um $k$-morfismo birracional.
Então $X$ é $k$-isomorfo a um $Y$-esquema obtido por explosão de um subesquema fechado $Y'$ de $Y$, não necessariamente reduzido.
\end{corollary}

\begin{proof}
\oldpage[III]{107}
Com efeito, $f$ é projetivo \sref[II]{II.5.5.5}[v] e, como $Y$ é projetivo sobre $k$, satisfaz-se a condição adicional de \sref{III.2.3.6}.
Basta, portanto, tomar o subesquema fechado $Y'$ de $Y$ definido pelo ideal coerente $\sh{J}$ de \sref{III.2.3.6}.
\end{proof}

\begin{env}[2.3.8]
\label{III.2.3.8}
\emph{Observação.}
No Capítulo~IV, ao estudar a noção de divisor, veremos que, se no enunciado de \sref{III.2.3.5} se supõe que os anéis locais $\sh{O}_y$, para $y\in Y$, são fatoriais---como acontece, por exemplo, quando $Y$ é não singular---, então $X$ pode ser obtido de $Y$ por explosão de um subpré-esquema fechado $Y'$ de $Y$ cujo espaço subjacente está contido em $Y-U$.
\end{env}


\subsection{Uma generalização do teorema fundamental}
\label{subsection:III.2.4}

\begin{theorem}[2.4.1]
\label{III.2.4.1}
Seja $Y$ um pré-esquema noetheriano, seja $\sh{S}$ uma $\sh{O}_Y$-álgebra quase-coerente de tipo finito, seja
\[
  f:X\longrightarrow Y
\]
um morfismo projetivo, ponhamos $\sh{S}'=f^*(\sh{S})$, e seja $\sh{M}$ um $\sh{S}'$-módulo quase-coerente de tipo finito.
Então:
\begin{enumerate}
  \item[(i) ] Para todo $p\in\bb Z$, $\RR^p f_*(\sh{M})$ é um $\sh{S}$-módulo de tipo finito.

  \item[(ii) ] Suponhamos ainda que $\sh{L}$ é um $\sh{O}_X$-módulo invertível amplo para $f$ e, para todo $n\in\bb Z$, ponhamos
  \[
    \sh{M}(n)=\sh{M}\otimes\sh{L}^{\otimes n}.
  \]
  Existe um inteiro $N$ tal que, para $n\geq N$,
  \[
  \label{III.2.4.1.1}
    \RR^p f_*(\sh{M}(n))=0
    \tag{2.4.1.1}
  \]
  para todo $p>0$, e o homomorfismo canônico
  \[
    f^*\bigl(f_*(\sh{M}(n))\bigr)\longrightarrow\sh{M}(n)
  \]
  \sref[0]{0.4.4.3} é sobrejetivo.
\end{enumerate}
\end{theorem}

\begin{proof}
Ponhamos
\[
  Y'=\Spec(\sh{S}),\qquad X'=\Spec(\sh{S}'),
\]
de modo que $X'=X\times_Y Y'$ \sref[II]{II.1.5.5}.
Sejam
\[
  g:Y'\longrightarrow Y,\qquad g':X'\longrightarrow X
\]
os morfismos estruturais, que são afins por definição, e seja
\[
  f'=f_{(Y')}:X'\longrightarrow Y'.
\]
Temos assim um diagrama comutativo
\[
  \xymatrix{
    X\ar[d]_f &
    X'\ar[l]_{g'}\ar[d]^{f'}\ar[dl]^{h}\\
    Y &
    Y'\ar[l]^{g},
  }
\]
onde $f'$ é projetivo \sref[II]{II.5.5.5}[iii] e
\[
  h=f\circ g'=g\circ f'.
\]

\emph{(i) } Seja $\widetilde{\sh{M}}$ o $\sh{O}_{X'}$-módulo associado ao $\sh{S}'$-módulo quase-coerente $\sh{M}$ quando $X'$ é considerado um $X$-esquema afim \sref[II]{II.1.4.3}; assim,
\[
  \sh{M}=g'_*(\widetilde{\sh{M}}).
\]
Como $\sh{M}$ é um $\sh{S}'$-módulo de tipo finito, $\widetilde{\sh{M}}$ é um $\sh{O}_{X'}$-módulo de tipo finito \sref[II]{II.1.4.5}.
Além disso, $h$ é de tipo finito, já que $g$ e $f'$ são de tipo finito \sref[II]{II.1.3.7} e \sref[I]{I.6.3.4}[ii].
Assim, $X'$ é noetheriano \sref[I]{I.6.3.7}, e $\widetilde{\sh{M}}$ é, portanto, coerente.

Como $g'$ é afim, o homomorfismo canônico
\[
  \RR^p f_*(\sh{M})
  \longrightarrow
  \RR^p h_*(\widetilde{\sh{M}})
\]
é bijetivo \sref{III.1.3.4}.
É, além disso, um homomorfismo de $\sh{S}$-módulos.
Com efeito, o homomorfismo canônico
\[
\label{III.2.4.1.2}
  g'_*(\sh{O}_{X'})
  \otimes_{\sh{O}_X}
  g'_*(\widetilde{\sh{M}})
  \longrightarrow
  g'_*(\widetilde{\sh{M}})
  \tag{2.4.1.2}
\]
define a estrutura de $\sh{S}'$-módulo sobre $\sh{M}$, já que
\[
  \sh{S}'=g'_*(\sh{O}_{X'}).
\]
Dá canonicamente lugar a um homomorfismo
\[
  f_*\bigl(g'_*(\sh{O}_{X'})\bigr)
  \otimes
  \RR^p f_*\bigl(g'_*(\widetilde{\sh{M}})\bigr)
  \longrightarrow
  \RR^p f_*\bigl(g'_*(\widetilde{\sh{M}})\bigr).
\]
\oldpage[III]{108}
Por \sref[0]{0.12.2.2}, e dado que \eqref{III.2.4.1.2} é obtido por sua vez, aplicando \sref[0]{0.4.2.2.1}, do homomorfismo
\[
  \sh{O}_{X'}\otimes\widetilde{\sh{M}}
  \longrightarrow\widetilde{\sh{M}}
\]
que define a estrutura de $\sh{O}_{X'}$-módulo sobre $\widetilde{\sh{M}}$, o diagrama
\[
  \xymatrix{
    f_*\bigl(g'_*(\sh{O}_{X'})\bigr)
    \otimes\RR^p f_*\bigl(g'_*(\widetilde{\sh{M}})\bigr)
    \ar[r]\ar[d] &
    \RR^p f_*\bigl(g'_*(\widetilde{\sh{M}})\bigr)\ar[d]\\
    h_*(\sh{O}_{X'})
    \otimes\RR^p h_*(\widetilde{\sh{M}})
    \ar[r] &
    \RR^p h_*(\widetilde{\sh{M}})
  }
\]
é comutativo \sref[0]{0.12.2.6}.
Compondo as flechas horizontais com o homomorfismo induzido pelo homomorfismo canônico
\[
  \sh{S}
  \longrightarrow f_*(f^*(\sh{S}))
  =f_*(\sh{S}')
  =f_*\bigl(g'_*(\sh{O}_{X'})\bigr)
  =h_*(\sh{O}_{X'}),
\]
obtemos a afirmação.

Por outro lado, como $g$ é afim e $f'$ é separado e quase compacto, o homomorfismo canônico
\[
  \RR^p h_*(\widetilde{\sh{M}})
  \longrightarrow
  g_*\bigl(\RR^p f'_*(\widetilde{\sh{M}})\bigr)
\]
é bijetivo \sref{III.1.4.14}.
Como antes, é um isomorfismo de $\sh{S}$-módulos, desta vez usando a comutatividade de \sref[0]{0.12.6.2}.
Agora, $f'$ é projetivo e $\widetilde{\sh{M}}$ é coerente, então $\RR^p f'_*(\widetilde{\sh{M}})$ é um $\sh{O}_{Y'}$-módulo coerente por \sref{III.2.2.1}.
Segue-se que
\[
  g_*\bigl(\RR^p f'_*(\widetilde{\sh{M}})\bigr)
\]
é um $\sh{S}$-módulo de tipo finito \sref[II]{II.1.4.5}.

\emph{(ii)} Ponhamos
\[
  \sh{L}'=g'^*(\sh{L}),
\]
um $\sh{O}_{X'}$-módulo invertível.
Para todo $n\in\bb Z$, salvo isomorfismo canônico,
\[
  g'_*\bigl(\widetilde{\sh{M}}\otimes\sh{L}'^{\otimes n}\bigr)
  =
  g'_*(\widetilde{\sh{M}})\otimes\sh{L}^{\otimes n}
  =
  \sh{M}(n)
\]
\sref[0]{0.5.4.10}.
Aplicando o raciocínio de \emph{(i) } a
\[
  \widetilde{\sh{M}}\otimes\sh{L}'^{\otimes n}
\]
mostra que
\[
  \RR^p f_*
  \bigl(g'_*(\widetilde{\sh{M}}\otimes\sh{L}'^{\otimes n})\bigr)
  \simeq
  g_*\bigl(\RR^p f'_*
  (\widetilde{\sh{M}}\otimes\sh{L}'^{\otimes n})\bigr).
\]
Agora, $\sh{L}'$ é amplo para $f'$ \sref[II]{II.4.6.13} [iii].
Deduz-se de \sref{III.2.2.1} que existe um inteiro $N$ tal que
\[
  \RR^p f'_*
  \bigl(\widetilde{\sh{M}}\otimes\sh{L}'^{\otimes n}\bigr)=0
\]
para todo $p>0$ e todo $n\geq N$.
Isso demonstra \eqref{III.2.4.1.1}.

Deduz-se mais uma vez de \sref{III.2.2.1} que pode ser escolhido $N$ de modo que, para $n\geq N$, o homomorfismo canônico
\[
  f'^*\bigl(f'_*(\widetilde{\sh{M}}\otimes\sh{L}'^{\otimes n})\bigr)
  \longrightarrow
  \widetilde{\sh{M}}\otimes\sh{L}'^{\otimes n}
\]
seja sobrejetivo.
Como $g'_*$ é exato \sref[II]{II.1.4.4}, o homomorfismo correspondente
\[
  g'_*\Bigl(
    f'^*\bigl(f'_*(\widetilde{\sh{M}}\otimes\sh{L}'^{\otimes n})\bigr)
  \Bigr)
  \longrightarrow
  g'_*(\widetilde{\sh{M}}\otimes\sh{L}'^{\otimes n})
  =\sh{M}(n)
\]
é sobrejetivo.
Mas \sref[II]{II.1.5.2} dá
\[
  g'_*\Bigl(
    f'^*\bigl(f'_*(\widetilde{\sh{M}}\otimes\sh{L}'^{\otimes n})\bigr)
  \Bigr)
  =
  f^*\Bigl(
    g_*\bigl(f'_*(\widetilde{\sh{M}}\otimes\sh{L}'^{\otimes n})\bigr)
  \Bigr),
\]
e $g_*\circ f'_*=f_*\circ g'_*$.
Consequentemente,
\[
  g'_*\Bigl(
    f'^*\bigl(f'_*(\widetilde{\sh{M}}\otimes\sh{L}'^{\otimes n})\bigr)
  \Bigr)
  =
  f^*\bigl(f_*(\sh{M}(n))\bigr),
\]
o que conclui a demonstração.
\end{proof}

\begin{env}[2.4.2]
\label{III.2.4.2}
Aplicaremos \sref{III.2.4.1}, em particular, quando $\sh{S}$ for uma $\sh{O}_Y$-álgebra graduada em graus positivos e
\[
  \sh{M}=\bigoplus_{k\in\bb Z}\sh{M}_k
\]
seja um $\sh{S}'$-módulo graduado.
\oldpage[III]{109}
Sob as mesmas hipóteses de finitude sobre $\sh{S}$ e $\sh{M}$, deduz-se de \sref{III.2.4.1} que
\[
  \bigoplus_{k\in\bb Z}\RR^p f_*(\sh{M}_k)
\]
é um $\sh{S}$-módulo de tipo finito para todo $p$.
Sob as hipóteses adicionais de \sref{III.2.4.1}[ii], existe também um inteiro $N$ tal que, para $n\geq N$,
\[
  \RR^p f_*(\sh{M}_k(n))=0
\]
para todo $p>0$ e todo $k\in\bb Z$, e o homomorfismo canônico
\[
  f^*\bigl(f_*(\sh{M}_k(n))\bigr)
  \longrightarrow
  \sh{M}_k(n)
\]
é sobrejetivo para todo $k\in\bb Z$.
\end{env}


\subsection{Característica de Euler--Poincar\'e e polinômio de Hilbert}
\label{subsection:III.2.5}

\begin{env}[2.5.1]
\label{III.2.5.1}
Seja $A$ um anel artiniano e seja $X$ um $A$-esquema projetivo sobre
\[
  Y=\Spec(A).
\]
Para todo $\sh{O}_X$-módulo coerente $\sh{F}$, os $A$-módulos
\[
  \HH^i(X,\sh{F})\qquad(i\geq0)
\]
são do tipo finito \sref{III.2.2.1} e, portanto, têm comprimento finito, já que $A$ é artiniano.
Além disso, \sref{III.2.2.1} mostra que $\HH^i(X,\sh{F})=0$ salvo para um número finito de $i\geq0$.
Assim, o inteiro
\[
\label{III.2.5.1.1}
  \chi_A(\sh{F})
  =
  \sum_{i=0}^{\infty}
  (-1)^i\operatorname{length}_A\bigl(\HH^i(X,\sh{F})\bigr)
  \tag{2.5.1.1}
\]
é definido para todo $\sh{O}_X$-módulo coerente $\sh{F}$.
Quando $A$ é um anel local artiniano, $\chi_A(\sh{F})$ chama-se a \emph{característica de Euler--Poincar\'e} de $\sh{F}$ em relação a $A$.
Para $\sh{F}=\sh{O}_X$, o inteiro $\chi_A(\sh{O}_X)$ chama-se o \emph{gênero aritmético} de $X$ em relação a $A$.
\end{env}

\begin{proposition}[2.5.2]
\label{III.2.5.2}
Seja
\[
  0\longrightarrow\sh{F}'
  \longrightarrow\sh{F}
  \longrightarrow\sh{F}''
  \longrightarrow0
\]
uma sequência exata de $\sh{O}_X$-módulos coerentes.
Então
\[
\label{III.2.5.2.1}
  \chi_A(\sh{F})
  =
  \chi_A(\sh{F}')+\chi_A(\sh{F}'').
  \tag{2.5.2.1}
\]
\end{proposition}

\begin{proof}
Os módulos de cohomologia de $\sh{F}$, $\sh{F}'$ e $\sh{F}''$ anulam-se em todos os graus suficientemente grandes.
Existe, portanto, um inteiro $r>0$ para o qual a longa sequência exata de cohomologia tem a forma
\[
  0\longrightarrow\HH^0(X,\sh{F}')
  \longrightarrow\HH^0(X,\sh{F})
  \longrightarrow\HH^0(X,\sh{F}'')
  \longrightarrow\HH^1(X,\sh{F}')
  \longrightarrow\cdots
\]
\[
  \cdots\longrightarrow\HH^r(X,\sh{F}')
  \longrightarrow\HH^r(X,\sh{F})
  \longrightarrow\HH^r(X,\sh{F}'')
  \longrightarrow0.
\]
Em uma sequência exata de $A$-módulos de comprimento finito com termos nulos em ambas as extremidades, a soma alternada dos comprimentos é nula \sref[0]{0.11.10.1}.
A aplicação deste fato dá imediatamente \eqref{III.2.5.2.1}.
\end{proof}

O resultado de \sref{III.2.5.2} aplica-se sempre que $X$ seja um $A$-esquema quase compacto e que os $A$-módulos $\HH^i(X,\sh{F})$ sejam de tipo finito para todo $\sh{O}_X$-módulo coerente $\sh{F}$ \sref{III.1.4.12}.

\begin{theorem}[2.5.3]
\label{III.2.5.3}
Seja $A$ um anel local artiniano, seja $X$ um esquema projetivo sobre
\[
  Y=\Spec(A),
\]
seja $\sh{L}$ um $\sh{O}_X$-módulo invertível muito amplo relativo a $Y$, e seja $\sh{F}$ um $\sh{O}_X$-módulo coerente.
Para todo $n\in\bb Z$, ponhamos
\[
  \sh{F}(n)=\sh{F}\otimes_{\sh{O}_X}\sh{L}^{\otimes n}.
\]
\begin{enumerate}
  \item[(i) ] Existe um único polinômio $P\in\bb Q[T]$ tal que
  \[
    \chi_A(\sh{F}(n))=P(n)
  \]
  para todo $n\in\bb Z$.
  O polinômio $P$ chama-se o \emph{polinômio de Hilbert} de $\sh{F}$ em relação a $A$.

  \item[(ii) ] Para todo $n$ suficientemente grande,
  \[
    \chi_A(\sh{F}(n))
    =
    \operatorname{length}_A\Gamma(X,\sh{F}(n)).
  \]

  \item[(iii)] O coeficiente do termo de maior grau de $\chi_A(\sh{F}(n))$ não é negativo.
\end{enumerate}
\end{theorem}

Além disso, demonstraremos no Capítulo~IV, na seção dedicada à dimensão, que \emph{o grau de $\chi_A(\sh{F}(n))$ é igual à dimensão do suporte de $\sh{F}$}.

\begin{proof}
Para todo $n$ suficientemente grande, tem-se
\[
  \HH^i(X,\sh{F}(n))=0
\]
para todo $i>0$ \sref{III.2.2.1}.
\oldpage[III]{110}
Consequentemente,
\[
  \chi_A(\sh{F}(n))
  =
  \operatorname{length}_A\HH^0(X,\sh{F}(n))
  =
  \operatorname{length}_A\Gamma(X,\sh{F}(n))
\]
para todo $n$ suficientemente grande, o que demonstra \emph{(ii) }.
Segue-se que $\chi_A(\sh{F}(n))\geq0$ para todo $n$ suficientemente grande e, portanto, \emph{(iii) } deduz-se de \emph{(i) }.
A afirmação de unicidade de \emph{(i) } é imediata, por isso resta provar a existência de $P$.

Primeiro, vamos mostrar que podemos supor
\[
  \mathfrak{m}\sh{F}=0,
\]
onde $\mathfrak{m}$ é o ideal maximal de $A$.
Com efeito, existe um inteiro $s>0$ tal que $\mathfrak{m}^s=0$ e, portanto, $\sh{F}(n)$ possui a filtração finita
\[
  \sh{F}(n)
  \supset\mathfrak{m}\sh{F}(n)
  \supset\cdots
  \supset\mathfrak{m}^{s-1}\sh{F}(n)
  \supset0.
\]
Uma indução que usa \eqref{III.2.5.2.1} dá
\[
  \chi_A(\sh{F}(n))
  =
  \sum_{k=1}^{s}
  \chi_A\left(
    \mathfrak{m}^{k-1}\sh{F}(n)/
    \mathfrak{m}^{k}\sh{F}(n)
  \right).
\]
Mas
\[
  \mathfrak{m}^{k-1}\sh{F}(n)/
  \mathfrak{m}^{k}\sh{F}(n)
  =
  \sh{F}'_k(n),
  \qquad
  \sh{F}'_k=
  \mathfrak{m}^{k-1}\sh{F}/\mathfrak{m}^{k}\sh{F},
\]
o que demonstra a redução.

Então, suponhamos que $\mathfrak{m}\sh{F}=0$.
Seja $X'$ o subesquema fechado de $X$ obtido como imagem inversa, pelo morfismo estrutural
\[
  X\longrightarrow\Spec(A),
\]
do único ponto fechado de $\Spec(A)$, e seja $j:X'\to X$ a injeção canônica.
Então
\[
  \sh{F}=j_*(\sh{F}')
\]
para um $\sh{O}_{X'}$-módulo coerente $\sh{F}'$, e $X'$ é projetivo sobre $\Spec(K)$, onde
\[
  K=A/\mathfrak{m}.
\]
Se $\sh{L}'=j^*(\sh{L})$, então $\sh{L}'$ é muito amplo relativo a $\Spec(K)$ \sref[II]{II.4.4.10}, e
\[
  \sh{F}(n)=j_*(\sh{F}'(n)),
  \qquad
  \sh{F}'(n)
  =
  \sh{F}'\otimes_{\sh{O}_{X'}}\sh{L}'^{\otimes n}
\]
\sref[0]{0.5.4.10}.
Segue-se que
\[
  \chi_A(\sh{F}(n))
  =
  \chi_K(\sh{F}'(n))
\]
(G, II, 4.9.1), e ficamos reduzidos ao caso em que $A$ é um corpo.

Agora, $X$ pode ser considerado um subesquema fechado de
\[
  P=\mathbf{P}_A^r
\]
para um $r$ conveniente \sref[II]{II.5.5.4}[ii].
Se $i:X\to P$ é a injeção canônica, o mesmo raciocínio dá
\[
  \chi_A(\sh{F}(n))
  =
  \chi_A\bigl(i_*(\sh{F})(n)\bigr).
\]
Podemos, portanto, limitar-nos ao caso
\[
  X=\mathbf{P}_A^r=\Proj(S),
  \qquad
  S=A[T_0,\ldots,T_r],
\]
onde $A$ é um corpo.

Neste caso, $\sh{F}=\widetilde{M}$ para um $S$-módulo graduado $M$ de tipo finito \sref[II]{II.2.7.8}.
Pelo teorema das sicigias de Hilbert (M, VIII, 6.5), $M$ possui uma resolução finita por $S$-módulos livres graduados de tipo finito:
\[
  0\longrightarrow L_q
  \longrightarrow L_{q-1}
  \longrightarrow\cdots
  \longrightarrow L_1
  \longrightarrow M
  \longrightarrow0.
\]
Como o funtor $M\mapsto\widetilde{M}$ é exato \sref[II]{II.2.5.4}, obtemos uma sequência exata
\[
  0\longrightarrow\widetilde{L}_q
  \longrightarrow\widetilde{L}_{q-1}
  \longrightarrow\cdots
  \longrightarrow\widetilde{L}_1
  \longrightarrow\widetilde{M}
  \longrightarrow0,
\]
e, portanto, para todo $n\in\bb Z$, uma sequência exata
\[
  0\longrightarrow\widetilde{L}_q(n)
  \longrightarrow\widetilde{L}_{q-1}(n)
  \longrightarrow\cdots
  \longrightarrow\widetilde{L}_1(n)
  \longrightarrow\widetilde{M}(n)
  \longrightarrow0.
\]
Uma indução sobre $q$, usando \sref{III.2.5.2}, dá
\[
  \chi_A(\widetilde{M}(n))
  =
  \sum_{j=1}^{q}
  (-1)^{j+1}\chi_A(\widetilde{L}_j(n)).
\]
Para provar \emph{(i) }, ficamos, portanto, reduzidos ao caso em que $M$ é graduado, livre e de tipo finito e, portanto, ao caso
\[
  M=S(h)
\]
para algum $h\in\bb Z$.
Então
\[
  \widetilde{M}(n)
  =
  \widetilde{M(n)}
  =
  \widetilde{S(n+h)}
\]
\sref[II]{II.2.5.15}, então o teorema deduz-se do seguinte lema.
\end{proof}

\oldpage[III]{111}
\begin{lemma}[2.5.3.1]
\label{III.2.5.3.1}
Seja $A$ um corpo, seja $r>0$ um inteiro e seja
\[
  X=\mathbf{P}_A^r.
\]
Então
\[
  \chi_A(\sh{O}_X(n))
  =
  \binom{n+r}{r}
\]
para todo $n\in\bb Z$.
\end{lemma}

\begin{proof}
Para $n\geq0$,
\[
  \chi_A(\sh{O}_X(n))
  =
  \operatorname{length}_A\HH^0(X,\sh{O}_X(n)),
\]
que é o número de monômios nos $T_i$ de grau total $n$, ou seja,
\[
  \binom{n+r}{r}
\]
\sref{III.2.1.12}.

Para $n\leq-r-1$, tem-se analogamente
\[
  \chi_A(\sh{O}_X(n))
  =
  (-1)^r\operatorname{length}_A\HH^r(X,\sh{O}_X(n)).
\]
Escrevendo $n=-r-h$, a dimensão de $\HH^r(X,\sh{O}_X(n))$ sobre $A$ é o número de sequências
\[
  (p_i)_{0\leq i\leq r}
\]
de inteiros positivos que satisfaçam
\[
  \sum_{i=0}^{r}p_i=r+h
\]
\sref{III.2.1.12}.
Equivalentemente, é o número de sequências de inteiros não negativos
\[
  (q_i)_{0\leq i\leq r}
\]
tais que
\[
  \sum_{i=0}^{r}q_i=h-1.
\]
Este número é
\[
  \binom{h+r-1}{r}
  =
  (-1)^r\binom{n+r}{r}.
\]
Finalmente, se $-r\leq n<0$, então
\[
  \binom{n+r}{r}=0,
\]
e $\HH^i(X,\sh{O}_X(n))=0$ para todo $i\geq0$ \sref{III.2.1.12}.
Isto demonstra o lema.
\end{proof}

\begin{corollary}[2.5.4]
\label{III.2.5.4}
Seja $A$ um anel local artiniano, seja $S$ uma $A$-álgebra graduada de tipo finito gerada por $S_1$, seja $M$ um $S$-módulo graduado de tipo finito e seja
\[
  X=\Proj(S).
\]
Então, para todo $n$ suficientemente grande,
\[
  \chi_A(\widetilde{M}(n))
  =
  \operatorname{length}_A(M_n).
\]
\end{corollary}

\begin{proof}
Para todo $n$ suficientemente grande, os módulos
\[
  M_n
  \quad\text{e}\quad
  \Gamma(X,\widetilde{M}(n))
\]
são isomorfos \sref{III.2.3.1}.
\end{proof}


\subsection{Aplicação: criterios de amplitud}
\label{subsection:III.2.6}

\begin{proposition}[2.6.1]
\label{III.2.6.1}
Seja $Y$ um pré-esquema noetheriano, seja $f:X\to Y$ um morfismo próprio,
e seja $\sh{L}$ um $\sh{O}_X$-módulo invertível.
As seguintes condições são equivalentes:
\begin{enumerate}
  \item[\emph{a) }] $\sh{L}$ é amplo para $f$.
  \item[\emph{b)}] Para todo $\sh{O}_X$-módulo coerente $\sh{F}$, existe
  um inteiro $N$ tal que, para $n\geq N$,
  \[
    \RR^qf_*(\sh{F}\otimes\sh{L}^{\otimes n})=0
  \]
  para todo $q>0$.
  \item[\emph{c)}] Para todo ideal coerente $\sh{J}$ de $\sh{O}_X$, existe
  um inteiro $N$ tal que, para $n\geq N$,
  \[
    \RR^1f_*(\sh{J}\otimes\sh{L}^{\otimes n})=0.
  \]
\end{enumerate}
\end{proposition}

\begin{proof}
Vimos que \emph{a) } implica \emph{b) } \sref{III.2.2.1},~(ii).
É imediato que \emph{b) } implica \emph{c) }, e resta demonstrar
que \emph{c) } implica \emph{a) }.
Podemos limitar-nos ao caso em que $Y$ é afim \sref[II]{II.4.6.4}.
Basta mostrar que, quando $h$ percorre as seções de
$\sh{L}^{\otimes n}$ sobre $X$, com $n>0$, os subconjuntos abertos afins $X_h$
formam um recobrimento de $X$ \sref[II]{II.4.5.2}.

Para isso, seja $x$ um ponto fechado de $X$ e seja $U$ uma vizinhança
aberta afim de $x$.
Encontraremos um inteiro $n$ e uma seção
\[
  h\in\Gamma(X,\sh{L}^{\otimes n})
\]
tal que $x\in X_h\subset U$.
O subconjunto aberto $X_h$ será então afim \sref[I]{I.1.3.6}.
A união de todos esses $X_h$ é um subconjunto aberto de $X$ que contém todo
ponto fechado e, portanto, é todo $X$, já que $X$ é noetheriano
\sref[I]{I.6.3.7} e \sref[0]{0.2.1.3}.

Sejam $\sh{J}$ e $\sh{J}'$, respectivamente, os ideais quase-coerentes de
$\sh{O}_X$ que definem os subesquemas fechados reduzidos cujos espaços
subjacentes são $X-U$ e $(X-U)\cup\{x\}$, respectivamente \sref[I]{I.5.2.1}.
Os ideais $\sh{J}$ e $\sh{J}'$ são coerentes \sref[I]{I.6.1.1},
$\sh{J}'\subset\sh{J}$, e
\[
  \sh{J}''=\sh{J}/\sh{J}'
\]
é um $\sh{O}_X$-módulo coerente \sref[0]{0.5.3.3}, com suporte
$\{x\}$, e $\sh{J}''_x=\kres(x)$.
Como $\sh{L}$ é localmente livre, a sequência
\[
  0\longrightarrow
  \sh{J}'\otimes\sh{L}^{\otimes n}
  \longrightarrow
  \sh{J}\otimes\sh{L}^{\otimes n}
  \longrightarrow
  \sh{J}''\otimes\sh{L}^{\otimes n}
  \longrightarrow0
\]
é exata para todo $n$.
Por hipótese, para todo $n$ suficientemente grande,
\[
  \HH^1(X,\sh{J}'\otimes\sh{L}^{\otimes n})=0.
\]

\oldpage[III]{112}

A sequência exata de cohomologia mostra, portanto, que
\[
  \Gamma(X,\sh{J}\otimes\sh{L}^{\otimes n})
  \longrightarrow
  \Gamma(X,\sh{J}''\otimes\sh{L}^{\otimes n})
\]
é sobrejetivo.
Uma seção $g$ de $\sh{J}''\otimes\sh{L}^{\otimes n}$ sobre $X$ tal que
$g(x)\neq0$ é, portanto, a imagem de uma seção
\[
  h\in\Gamma(X,\sh{J}\otimes\sh{L}^{\otimes n})
  \subset\Gamma(X,\sh{L}^{\otimes n});
\]
a inclusão deduz-se de \sref[0]{0.5.4.1}.
Por construção, $h(x)\neq0$, enquanto $h(z)=0$ para todo $z\notin U$.
Assim, $x\in X_h\subset U$, o que prova a afirmação.
\end{proof}

\begin{proposition}[2.6.2]
\label{III.2.6.2}
Seja $Y$ um pré-esquema noetheriano, seja $f:X\to Y$ um morfismo de
tipo finito, seja $g:X'\to X$ um morfismo finito sobrejetivo, seja
$\sh{L}$ um $\sh{O}_X$-módulo invertível, e ponhamos
\[
  \sh{L}'=g^*(\sh{L}).
\]
Suponhamos que existe um subconjunto $Z$ de $X$, próprio sobre $Y$
\sref[II]{II.5.4.10}, tal que, para todo $x\in X-Z$, ou $X$ é
normal em $x$, ou
\[
  \bigl(g_*(\sh{O}_{X'})\bigr)_x
\]
é um $\sh{O}_{X,x}$-módulo livre.
Então $\sh{L}$ é amplo para $f$ se e apenas se $\sh{L}'$ é amplo para
$f\circ g$.
\end{proposition}

\begin{proof}
\begin{env}[2.6.2.1]
\label{III.2.6.2.1}
Como $g$ é afim, a condição é necessária \sref[II]{II.5.1.12}.
Para provar que é suficiente, podemos supor que $Y$ é afim.
\sref[II]{II.4.6.4}.
Primeiro vamos mostrar que também podemos limitar-nos ao caso em que $X$ é
reducido.

Seja $j:X_{\mathrm{red}}\to X$ a imersão canônica, e ponhamos
\[
  X_1=X_{\mathrm{red}},
  \qquad
  X'_1=X'\times_X X_1.
\]
Então temos o diagrama comutativo
\[
  \xymatrix{
    X'\ar[d]_g &
    X'_1\ar[l]_{j'}\ar[d]^{g_1}\\
    X &
    X_1\ar[l]^{j}
  }
  \label{III.2.6.2.2}\tag{2.6.2.2}
\]
O morfismo $f\circ j$ é do tipo finito \sref[I]{I.6.3.4}, e
$g_1$ é finito \sref[II]{II.6.1.5},~(iii).
Se $\sh{L}'$ é amplo para $f\circ g$, então
$j'^*(\sh{L}')$ é amplo para $f\circ g\circ j'$, porque $j'$ é uma
imersão fechada \sref[II]{II.5.1.12} e \sref[I]{I.4.3.2}.
Se $Z_1=j^{-1}(Z)$, então $Z_1$ é próprio sobre $Y$
\sref[II]{II.5.4.10}.
Além disso, se $X$ é normal em um ponto $x$, também o é $X_{\mathrm{red}}$
nesse ponto; e, se
$\bigl(g_*(\sh{O}_{X'})\bigr)_x$ é um $\sh{O}_{X,x}$-módulo livre,
então \sref[II]{II.1.5.2} mostra que
\[
  \bigl((g_1)_*(\sh{O}_{X'_1})\bigr)_x
\]
é um $\sh{O}_{X_1,x}$-módulo livre.
Finalmente, como $X$ é noetheriano \sref[I]{I.6.3.7}, a amplitude de
$j^*(\sh{L})$ implica a de $\sh{L}$ \sref[II]{II.4.5.14}.
Como
\[
  j'^*(\sh{L}')=g_1^*(j^*(\sh{L})),
\]
Isto demonstra a redução anunciada.
Em seguida, assumimos que $Y$ é afim e que $X$ é reduzido.
\end{env}

As hipóteses de \sref[II]{II.6.6.11} são agora satisfeitas.
Existem um $Y$-pré-esquema reduzido $X_2$ e um
$Y$-morfismo finito birracional
\[
  h:X_2\to X
\]
cuja restrição a $h^{-1}(X-Z)$ é um isomorfismo sobre $X-Z$, e
tal que $h^*(\sh{L})$ é amplo.
Substituindo $X'$ por $X_2$, ficamos reduzidos a demonstrar a proposição sob
a hipótese adicional de que $g$ possui as propriedades que acabamos de listar para
$h$.
Seguimos denotando por $Z$ um subesquema fechado de $X$, próprio sobre
$Y$ \sref[II]{II.5.4.10}, cujo espaço subjacente é o subconjunto dado
$Z$.

\begin{env}[2.6.2.3]
\label{III.2.6.2.3}
Seja $X_1$ um subpreesquema fechado de $X$, seja $j:X_1\to X$ a
imersão canônica, e seja
\[
  X'_1=g^{-1}(X_1)=X'\times_X X_1,
\]
com imersão canônica $j':X'_1\to X'$.
O quadrado correspondente é novamente \eqref{III.2.6.2.2}.
Ponhamos
\[
  \sh{L}_1=j^*(\sh{L}),
  \qquad
  \sh{L}'_1=j'^*(\sh{L}')=g_1^*(\sh{L}_1).
\]
Então $\sh{L}'_1$ é amplo para $f\circ g\circ j'$
\sref[II]{II.5.1.12}.
Se $Z_1=j^{-1}(Z)$, o subpré-esquema fechado $Z_1$ de $X_1$ é próprio
sobre $Y$ \sref[II]{II.5.4.2},~(ii).
Em outras palavras, as hipóteses de \sref{III.2.6.2} verificam-se para
$X_1$, $\sh{L}_1$, $g_1$ e $Z_1$.
\end{env}

\oldpage[III]{113}

Isso nos permite demonstrar \sref{III.2.6.2} por indução noetheriana
\sref[0]{0.2.2.2} Quando a restrição de $g$ a $g^{-1}(X-Z)$ é um
isomorfismo sobre $X-Z$; como foi visto em \sref{III.2.6.2.1}, esse caso é
suficiente.
Basta mostrar que, se a conclusão de \sref{III.2.6.2} valer
para $\sh{L}_1$ sobre todo subpré-esquema fechado $X_1$ de $X$ cujo
o espaço subjacente não é $X$, então também vale para $\sh{L}$.

\begin{env}[2.6.2.4]
\label{III.2.6.2.4}
Ponhamos
\[
  \sh{A}=\sh{O}_X,
  \qquad
  \sh{B}=g_*(\sh{O}_{X'}).
\]
Como $g$ é birracional e $X$ é reduzido, $\sh{B}$ é uma
sub-$\sh{A}$-álgebra do feixe $\sh{R}(X)$ de anéis totais de frações,
e é um $\sh{A}$-módulo coerente.
Além disso,
\[
  \sh{B}|_{X-Z}=\sh{A}|_{X-Z}.
\]
Seja $\sh{K}$ o condutor de $\sh{B}$ sobre $\sh{A}$, ou seja, o maior
sub-$\sh{A}$-módulo quase-coerente de $\sh{A}$ tal que
\[
  \sh{B}\sh{K}\subset\sh{A}.
\]
Equivalentemente, $\sh{K}$ é o anulador do
$\sh{A}$-módulo $\sh{B}/\sh{A}$ \sref[0]{0.5.3.7}.
Segue-se que
\[
  \sh{B}\sh{K}=\sh{K}.
\]
É claro que $\sh{K}_x=\sh{A}_x$ em todo ponto que admita uma vizinhança
$W_x$ sobre o qual $g:g^{-1}(W_x)\to W_x$ é um isomorfismo.
Isto ocorre, em particular, em todo ponto de $X-Z$, e em uma
vizinhança de todo ponto genérico de uma componente irredutível de
$X$.

Seja
\[
  Z_1=\Spec(\sh{A}/\sh{K})
\]
o subpré-esquema fechado de $X$ definido por $\sh{K}$.
É próprio sobre $Y$, porque seu espaço subjacente é fechado em $Z$
\sref[II]{II.5.4.10}.
A definição de $\sh{K}$ também dá
\[
  \sh{B}|_{X-Z_1}=\sh{A}|_{X-Z_1}.
\]
Podemos, portanto, reduzir-nos ao caso
\[
  Z=\Spec(\sh{A}/\sh{K}).
\]
Como $X-Z_1$ é um subconjunto aberto não vazio de $X$, podemos também
supor que o espaço subjacente de $Z$ seja diferente de $X$.
\end{env}

\begin{env}[2.6.2.5]
\label{III.2.6.2.5}
Consideremos $X'$ como $\Spec(\sh{B})$, já que $g$ é afim, e seja
\[
  \sh{K}'=\widetilde{\sh{K}}.
\]
Este é um ideal coerente de $\sh{O}_{X'}$ que satisfaz
$g_*(\sh{K}')=\sh{K}$ \sref[II]{II.1.4.1}.
O subpré-esquema fechado
\[
  Z'=g^{-1}(Z)=Z\times_X X'
\]
de $X'$ é definido por $\sh{K}'$ e é igual a
\[
  \Spec(\sh{B}/\sh{K})
\]
\sref[II]{II.1.4.10}.
O morfismo $h:Z'\to Z$ é finito \sref[II]{II.6.1.5},~(iii);
portanto, $Z'$ é próprio sobre $Y$ \sref[II]{II.6.1.11} e
\sref[II]{II.5.4.2},~(ii).

Devemos demonstrar que, para todo $x\in X$ e toda vizinhança aberta $U$
de $x$, existem um inteiro $n>0$ e uma seção $s$ de
$\sh{L}^{\otimes n}$ sobre $X$ tais que
\[
  x\in X_s\subset U
\]
\sref[II]{II.4.5.2}.
Distinguimos dos casos.

\smallskip
\noindent
$1^\circ$ Vamos primeiro supor que $x\in X-Z$.
Podemos, evidentemente, supor que $U\subset X-Z$, de modo que
$U'=g^{-1}(U)$ não intersecta $Z'$.
Como $\sh{L}'$ é amplo, existem $n>0$ e uma seção $s'$ de
$\sh{L}'^{\otimes n}$ sobre $X'$ tais que
\[
  x'=g^{-1}(x)\in X'_{s'}\subset g^{-1}(U)
\]
\sref[II]{II.4.5.2}.
Também podemos supor que
$\sh{K}'\otimes\sh{L}'^{\otimes n}$ é gerado por suas seções globais
\sref[II]{II.4.5.5}.
Como $\sh{K}'_{x'}=\sh{O}_{X',x'}$, uma destas seções, digamos $s''$, é
não nula em $x'$.
Multiplicando por $s'$ e substituindo assim $n$ por $2n$, podemos obter
que
\[
  x'\in X'_{s''}\subset g^{-1}(U).
\]
Por \sref[0]{0.5.4.10}, existe um isomorfismo canônico
\[
  \Gamma(X,\sh{K}\otimes\sh{L}^{\otimes n})
  \xrightarrow{\sim}
  \Gamma(X',\sh{K}'\otimes\sh{L}'^{\otimes n}).
\]
A seção $s$ de $\sh{K}\otimes\sh{L}^{\otimes n}$ correspondente a
$s''$ possui manifestamente as propriedades exigidas.

\smallskip
\noindent
$2^\circ$ Agora vamos supor que $x\in Z$.
Seja $\sh{J}$ o ideal coerente de $\sh{O}_X$ que define o
subpreesquema fechado reduzido cujo espaço subjacente é $X-U$, e consideremos em
$\sh{B}$ os ideais coerentes

\oldpage[III]{114}

\[
  \sh{J}\sh{B}
  \qquad\text{e}\qquad
  \sh{J}\sh{K}
  =
  \sh{J}(\sh{K}\sh{B})
  =
  \sh{K}(\sh{J}\sh{B}).
\]
Temos o diagrama de inclusões
\[
  \xymatrix@R=1.5em{
    \sh{J}\sh{B}\ar@{^{(}->}[r] &
    \sh{B}\\
    \sh{J}\ar@{^{(}->}[r]\ar@{^{(}->}[u] &
    \sh{A}\ar@{^{(}->}[u]\\
    \sh{J}\sh{K}\sh{B}=\sh{J}\sh{K}
      \ar@{^{(}->}[r]\ar@{^{(}->}[u] &
    \sh{K}\ar@{^{(}->}[u]
  }
  \label{III.2.6.2.6}\tag{2.6.2.6}
\]
Seja $\sh{J}'=(\sh{J}\sh{B})^\sim$ o ideal coerente de
$\sh{O}_{X'}$ tal que
\[
  g_*(\sh{J}')=\sh{J}\sh{B}.
\]
Então
\[
  \sh{J}'\sh{K}'=(\sh{J}\sh{K}\sh{B})^\sim,
  \qquad
  \sh{J}'/\sh{J}'\sh{K}'
  =
  (\sh{J}\sh{B}/\sh{J}\sh{K}\sh{B})^\sim
\]
\sref[II]{II.1.4.4}.
Como $\sh{J}|_V=\sh{J}\sh{K}|_V$ sobre todo subconjunto aberto $V$ disjunto
de $Z$, o suporte de $\sh{J}'/\sh{J}'\sh{K}'$ está contido em
$Z'$.
Como $Z'$ é próprio sobre $Y$, \sref{III.2.2.4} mostra que, para todo
$n$ suficientemente grande, a aplicação canônica
\[
  \Gamma(X',\sh{J}'\otimes\sh{L}'^{\otimes n})
  \longrightarrow
  \Gamma\bigl(
    X',
    (\sh{J}'/\sh{J}'\sh{K}')\otimes\sh{L}'^{\otimes n}
  \bigr)
\]
é sobrejetiva.
Por \sref[0]{0.5.4.10}, segue-se que a aplicação canônica
\[
  \Gamma(X,\sh{J}\sh{B}\otimes\sh{L}^{\otimes n})
  \longrightarrow
  \Gamma\bigl(
    X,
    (\sh{J}\sh{B}/\sh{J}\sh{K}\sh{B})
      \otimes\sh{L}^{\otimes n}
  \bigr)
\]
é sobrejetiva.

Sejam $i:Z\to X$ e $i':Z'\to X'$ as imersões canônicas.
São parte do diagrama comutativo
\[
  \xymatrix{
    X'\ar[d]_g &
    Z'\ar[l]_{i'}\ar[d]^h\\
    X &
    Z\ar[l]^{i}
  }
\]
Ponhamos
\[
  \sh{M}=i^*(\sh{L}),
  \qquad
  \sh{M}'=i'^*(\sh{L}').
\]
Como $\sh{L}'$ é amplo, $\sh{M}'$ é amplo.
\sref[II]{II.5.1.12}, e
\[
  \sh{M}'=h^*(\sh{M}).
\]
A hipótese de indução noetheriana, aplicada pois $Z\neq X$, mostra agora
que $\sh{M}$ é amplo.
Por conseguinte, para todo $n$ suficientemente grande,
\[
  i^*(\sh{J}/\sh{J}\sh{K})\otimes\sh{M}^{\otimes n}
\]
é gerado por suas seções sobre $Z$ \sref[II]{II.4.5.5}.
Como
\[
  \sh{J}/\sh{J}\sh{K}
  =
  i_*\bigl(i^*(\sh{J}/\sh{J}\sh{K})\bigr),
\]
\sref[0]{0.5.4.10} dá uma seção $s$ de
\[
  (\sh{J}/\sh{J}\sh{K})\otimes\sh{L}^{\otimes n}
\]
sobre $X$, para algum $n$ suficientemente grande, tal que $s(x)\neq0$.
De fato, $\sh{J}_x=\sh{A}_x$ pela definição de $\sh{J}$, enquanto
$\sh{K}_x\neq\sh{A}_x$ por hipótese.

O diagrama \eqref{III.2.6.2.6} mostra que $s$ é também uma seção de
\[
  (\sh{J}\sh{B}/\sh{J}\sh{K}\sh{B})
  \otimes\sh{L}^{\otimes n}
\]
sobre $X$.
É, portanto, a imagem canônica de uma seção
\[
  t\in\Gamma(X,\sh{J}\sh{B}\otimes\sh{L}^{\otimes n}).
\]
Por definição, a imagem de $t$ módulo
$\sh{J}\sh{K}\sh{B}\otimes\sh{L}^{\otimes n}$ pertence a
\[
  \frac{\sh{J}\otimes\sh{L}^{\otimes n}}
       {\sh{J}\sh{K}\sh{B}\otimes\sh{L}^{\otimes n}}.
\]
O diagrama \eqref{III.2.6.2.6} implica, portanto, que $t$ é uma seção
de $\sh{J}\otimes\sh{L}^{\otimes n}$ e, em particular, uma seção de
$\sh{L}^{\otimes n}$.
Temos $t(x)\neq0$, logo $x\in X_t$.
Pela definição de $\sh{J}$, a seção $t$ anula-se sobre $X-U$, o
suporte de $\sh{O}_X/\sh{J}$.
Assim, $X_t\subset U$, e a demonstração está terminada.
\end{env}
\end{proof}

\oldpage[III]{115}

\begin{env}[2.6.3]
\label{III.2.6.3}
\emph{Observação.}
Quando $X$ é próprio sobre $Y$, a Proposição~\sref{III.2.6.2} pode ser demonstrada
mais simplesmente raciocinando como no teorema de Chevalley
\sref[II]{II.6.7.1}, com a ajuda da Proposição~\sref{III.2.6.1} e
do lema \sref[II]{II.6.7.1.1}.
\end{env}
